.mt4
.mb6
.po4
.ll80
.ig
CHAPTER 32
D M S	S C R E E N    F O R M A T S
.sp2
.li
		Y O U R    C O M P A N Y    N A M E
			   PAYROLL SYSTEM
  00/00/00		PAYROLL SYSTEM MENU

   1.  PAY TRANSACTION ENTRY		15.  PRINT QUARTERLY 941 REPORT
   2.  PAY TRANSACTION PROOF		16.  END THE CURRENT QUARTER
   3.  PAY TRANSACTION SORT		17.  PRINT W2 REPORT
   4.  APPLY PAYROLL			18.  PRINT YEARLY 940 REPORT
   5.  PRINT PAYROLL JOURNAL		19.  END THE CURRENT YEAR
   6.  PRINT PAYCHECKS
   7.  PRINT CHECK REGISTER		20.  PRINT EMPLOYEE MASTER LIST
   8.  SET SYSTEM DATE			21.  PRINT HISTORY REPORT
					22.  PRINT VACATION REPORT
   9.  PARAMETER ENTRY			23.  PRINT ACCOUNT LIST
  10.  QUICK PARAMETER ENTRY
  11.  ACCOUNT ENTRY			24.  EMPLOYEE ENTRY
  12.  DEDUCTION/EARNING ENTRY		25.  EMPLOYEE DED/EARN AND COST DIST.
  13.  CREATE SORT PARAMETER FILES	26.  EMPLOYEE PAY RATE ENTRY
  14.  HISTORY ENTRY			27.  USER DEFINED PROGRAM

		    ENTER A SELECTION NUMBER [	]
		    ENTER TODAY'S DATE [        ]
.sp3
.he Appendix: Chapter 32
.li
		  Y O U R    C O M P A N Y    N A M E
			    PAYROLL SYSTEM
  00/00/00		    PARAMETER ENTRY

 ---------------------------------------------------------------------------
 PAY INTERVAL USED(S,M,W,B  IF MORE THAN ONE, SELECT SCREEN 2)[ 	   ]
 DATE FORM (1=MM/DD/YY, 2=DD/MM/YY)    [ ]	   LAST QUARTER ENDED	 [ ]
 ENDING DAY OF LAST PAY PERIOD	[	 ]		CURRENT YEAR	[  ]
 ---------------------------------------------------------------------------
 COMPANY
    NAME    [								 ]
    ADDRESS 1 [ 			     ]	CITY  [ 		 ]
    ADDRESS 2 [ 			     ]	STATE [  ]    ZIP  [	 ]
    ADDRESS 3 [ 			     ]	PHONE	   [		 ]
    GOVT ID	   FEDERAL		  STATE 	       LOCAL
    NUMBERS   [ 	      ]     [		    ]	 [		 ]
    ----------------------------------------------------------------------
    LINES PER PAGE  [  ]    SUPPRESS CONSOLE BELL [ ]	 DEBUGGING AID [ ]
.bp
.sp8
.li
		1. COMPANY ID AND GENERAL PARAMETERS
		2. PAY INTERVALS AND DEFAULT EMPLOYEE PAYROLL DATA
		3. PAYROLL SYSTEM CONFIGURATION AND GL INTERFACE

			    SELECT A SCREEN  [ ]

			ENTER TODAY'S DATE [        ]

		   ENTER THE SYSTEM FILE DRIVE	(A OR B) [ ]
.sp4
.li
		   Y O U R   C O M P A N Y    N A M E
			    PAYROLL SYSTEM
  00/00/00		    PARAMETER ENTRY

     -------------------------------------------------------------------
     |		     P A Y   I N T E R V A L   U S A G E	       |
     | SEMI-MONTHLY [ ] END OF LAST SEMI-MONTHLY PAY PERIOD [	     ] |
     | MONTHLY	    [ ] END OF LAST MONTHLY PAY PERIOD	    [	     ] |
     | BIWEEKLY     [ ] END OF LAST BIWEEKLY PAY PERIOD     [	     ] |
     | WEEKLY	    [ ] END OF LAST WEEKLY PAY PERIOD	    [	     ] |
     -------------------------------------------------------------------
	      D E F A U L T   E M P L O Y E E	V A L U E S
   -- RATE NAME  -----------  DEFAULT  ------  DESCRIPTION  ---------
    VACATION	     RATE   [	      ] (UNITS EARNED PER PAY PERIOD)
    SICK LEAVE	     RATE   [	      ] (UNITS EARNED PER PAY PERIOD)
 1 [		   ] RATE   [	      ] (DOLLARS PER UNIT)
 2 [		   ] FACTOR [	      ] (RATE 1 * FACTOR = DEFAULT RATE 2)
 3 [		   ] FACTOR [	      ] (RATE 1 * FACTOR = DEFAULT RATE 3)
 4 [		   ] EXCLUSION TYPE [ ] (RATE 4 CAN BE EXEMPTED FROM TAXES)
 --------------------------------------------------------------------------
  DEFAULT PAY INTERVAL (ANY USED INTERVAL) [ ]	MAX PAY FACTOR [	 ]
  DEFAULT PAY TYPE (H=HOURLY, S=SALARIED)  [ ]	NORMAL UNITS   [	 ]
.bp
.li
		    Y O U R   C O M P A N Y   N A M E
			    PAYROLL SYSTEM
  00/00/00		    PARAMETER ENTRY

		   D I S K   D R I V E	 A L L O C A T I O N
    SINGLE RECORD FILES 	  STATIC  FILES 	   PAY CYCLE FILES
    TAX TABLES	    [ ]        PAYROLL ACCOUNTS [ ]	   HOURS BATCH	[ ]
    SYSTEM DED/EARN [ ]        EMPLOYEE MASTER	[ ]	   PAY CHECK	[ ]
    COMPANY HISTORY [ ]        EMPLOYEE HISTORY [ ]	   PAY DETAIL	[ ]
    -----------------------------------------------------------------------
      G E N E R A L   L E D G E R      |      C H E C K   W R I T I N G
    GL INTERFACE USED (Y OR N) [ ]     |  COMPUTER CHECK PRINT (Y OR N) [ ]
    GL BATCH FILE OUTPUT DRIVE [ ]     |  VOID CHECKS OVER MAX (Y OR N) [ ]
    GL COA FILE SUFFIX	      [  ]     |  MAXIMUM AMOUNT     [		  ]
    GL DATA DRIVE	       [ ]     |-----------------------------------
    ------------------------------------------|    PAYROLL SYSTEM ACCOUNTS
     U S E R   D E F I N E D   P R O G R A M  |  OFFSET CASH  ACCOUNT [   ]
    PROGRAM USED [ ]  PROGRAM NAME [	    ] |  DEFAULT DIST ACCOUNT [   ]
    DESC.    [				    ] |
    -----------------------------------------------------------------------
	  DISTRIBUTE EMPLOYEE COST TO  1-5 GL ACCOUNTS (Y OR N) [ ]
	       (IF N, ONE ACCOUNT CAN BE ENTERED PER EMPLOYEE)
.sp5
.li
		Y O U R   C O M P A N Y   N A M E
			 PAYROLL SYSTEM
  00/00/00		  ACCOUNT ENTRY

	RECORD NO. [   ]

	   ACCOUNT NUMBER:[	 ]
	   ACCOUNT NAME:  [				 ]


	PLACE GL CHART OF ACCOUNTS FILE IN DRIVE: d
	REPLACE DISKETTES IN ORIGINAL CONFIGURATION

	TYPE RETURN TO CONTINUE:
.bp
.li
		      Y O U R	C O M P A N Y	 N A M E
			       PAYROLL SYSTEM
  00/00/00		   SYSTEMWIDE DEDUCTION MENU



		  1.  STANDARD DEDUCTION RATES
		  2.  FEDERAL WITHOLDING TAX TABLE ENTRY
		  3.  SYSTEM EARNING AND DEDUCTION INFORMATION



			   ENTER SELECTION NUMBER: [ ]
.sp6
.li
		  Y O U R    C O M P A N Y    N A M E
			     PAYROLL SYSTEM
  00/00/00		STANDARD DEDUCTION RATES

			      ACCT
		       USED    NO      RATE	LIMIT

	       FWT:	[ ]   [   ]
	       SWT:	[ ]   [   ]
	       LWT:	[ ]   [   ]
	       FICA:	[ ]   [   ]   [     ] [ 	  ]
	       CO FICA: [ ]   [   ]   [     ] [ 	  ]
	       SDI:	[ ]   [   ]   [     ] [ 	  ]
	       CO FUTA: [ ]   [   ]   [     ] [ 	  ]
	       FUTA MAX STATE CREDIT: [     ]
	       CO SUI:	[ ]   [   ]   [     ] [ 	  ]
	       EIC:	[ ]   [   ]   [     ] [ 	  ]
	       EIC EXCESS:	      [     ] [ 	  ]
.bp
.li
		       Y O U R	 C O M P A N Y	 N A M E
				PAYROLL SYSTEM
  00/00/00	      FEDERAL WITHHOLDING TAX TABLE ENTRY


		       ALLOWANCE AMOUNT:  [	      ]

		     M A R R I E D		  S I N G L E
		WAGES OVER:   PERCENT	   WAGES OVER:	 PERCENT
	   1.  [	   ]  [     ]	  [	      ]  [     ]
	   2.  [	   ]  [     ]	  [	      ]  [     ]
	   3.  [	   ]  [     ]	  [	      ]  [     ]
	   4.  [	   ]  [     ]	  [	      ]  [     ]
	   5.  [	   ]  [     ]	  [	      ]  [     ]
	   6.  [	   ]  [     ]	  [	      ]  [     ]
	   7.  [	   ]  [     ]	  [	      ]  [     ]

	       (Use ANNUAL tables from IRS CIRCULAR 'E' only!)
.sp7
.li
		    Y O U R    C O M P A N Y	N A M E
			      PAYROLL SYSTEM
  00/00/00	  SYSTEM EARNING AND DEDUCTION INFORMATION

USED  EARN/DED DESC   E/D ACCT. A/P   FACTOR   LIMITED	 LIMIT	  XLCD CAT
1[ ][		    ] [ ] [   ] [ ] [	       ] [ ] [		  ][ ][  ]
2[ ][		    ] [ ] [   ] [ ] [	       ] [ ] [		  ][ ][  ]
3[ ][		    ] [ ] [   ] [ ] [	       ] [ ] [		  ][ ][  ]
4[ ][		    ] [ ] [   ] [ ] [	       ] [ ] [		  ][ ][  ]
5[ ][		    ] [ ] [   ] [ ] [	       ] [ ] [		  ][ ][  ]
.bp
.li
		   Y O U R   C O M P A N Y   N A M E
			     PAYROLL SYSTEM
  00/00/00		     EMPLOYEE ENTRY

 EMPLOYEE NUMBER [     ] ------------------------------------------------
    |	  NAME		 [				]		|
    |	  ADDRESS     1  [				]		|
    |		      2  [				]		|
    |		  CITY	 [		    ]  STATE [	]  ZIP [     ]	|
    |		  PHONE  [	       ]				|
    ---------------------------------------------------------------------
    SOC SEC NO.   [	      ] 		 |    PENSION (Y/N)   [ ]
    BANK ACCT NO. [			    ]	 |    JOB CODE	    [	]
    BIRTH DATE	  [	   ]			 |    TAXING STATE   [	]
    SEX 	  [ ]				 |
    ---------------------------------------------------------------------
    START DATE [	]      |      EMPLOYMENT STATUS      [		]
    TERM. DATE [	]      |      (A=ACTIVE, L=ON LEAVE, T=TERMINATED)

	     THE EMPLOYEE MASTER AND HISTORY FILES ARE NOT PRESENT
	     TYPE  ""RUN"" TO CREATE BOTH FILES.    PRESS THE ESCAPE
	     KEY TO RETURN TO THE PAYROLL SYSTEM MENU [   ]
.sp4
.li
		 Y O U R   C O M P A N Y   N A M E
			  PAYROLL SYSTEM
  00/00/00	     EMPLOYEE EARN/DED AND COST

EMPL NO. [     ]    NAME				   SSN
 --------------------------------------------------------------------------
 C O S T   D I S T R I B U T I O N   |	  E X E M P T I O N S	F R O M
 O F   E M P L O Y E E	 G R O S S   |	   T A X    W I T H O L D I N G
 A N D	 E M P L O Y E R   C O S T   |	FICA EXEMPT [ ]  FEDERAL EXEMPT [ ]
 ACCOUNT [   ]	PERCENT  [	]    |	FUTA EXEMPT [ ]  STATE	 EXEMPT [ ]
 ACCOUNT [   ]	PERCENT  [	]    |	SUI  EXEMPT [ ]  LOCAL	 EXEMPT [ ]
 ACCOUNT [   ]	PERCENT  [	]    |	SDI  EXEMPT [ ]
 ACCOUNT [   ]	PERCENT  [	]    |	---------------------------------
 ACCOUNT [   ]	PERCENT  [	]    |	EXEMPTIONS FROM SYSTEM DEDUCTIONS
		TOTAL	  nnn.nn     |	   1[ ]  2[ ]  3[ ]  4[ ]  5[ ]
--------------------------------------------------------------------------
E M P L O Y E E   S P E C I F I C   D E D U C T I O N S / E A R N I N G S

USED  EARN/DED DESC   E/D ACCT. A/P   FACTOR   LIMITED	  LIMIT   XCLD CAT
1[ ][		    ] [ ] [   ] [ ] [	       ] [ ] [		  ][ ][  ]
2[ ][		    ] [ ] [   ] [ ] [	       ] [ ] [		  ][ ][  ]
3[ ][		    ] [ ] [   ] [ ] [	       ] [ ] [		  ][ ][  ]

			    ACCOUNT NUMBER [   ]

.bp
.li
		  Y O U R   C O M P A N Y   N A M E
			   PAYROLL SYSTEM
  00/00/00		 EMPLOYEE RATE ENTRY

 EMPL NO [     ]   NAME 				   SSN
---------------------------------------------------------------------------
   PAY INTERVAL  AND  PAY RATES    |	    T A X   W I T H O L D I N G
PAY TYPE (H=HOURLY,S=SALARIED) [ ] |		A L L O W A N C E S
PAY INTERVAL (S=SEMI-MONTHLY,  [ ] |  MARITAL STATUS(M=MARRIED,S=SINGLE)[ ]
(M=MONTHLY, W=WEEKLY, B=BIWEEKLY)  |	      FEDERAL  WH  ALLOWANCES  [  ]
rate1 desc. name  RATE [	 ] |	      STATE    WH  ALLOWANCES  [  ]
rate2 desc. name  RATE [	 ] |	      LOCAL    WH  ALLOWANCES  [  ]
rate3 desc. name  RATE [	 ] |  EARNED INCOME CREDIT USED (Y OR N)[ ]
rate4 desc. name  RATE [	 ] |  S T A T E   O F	C A L I F O R N I A
RATE 4 TAX EXCLUSION TYPE(1-4) [ ] |  S P E C I A L   I N F O R M A T I O N
AUTOGEN UNITS	       [	 ] |  HEAD OF HOUSEHOLD 	(Y OR N)[ ]
NORMAL UNITS	       [	 ] |  SPECIAL  WITHOLDING ALLOWANCES   [  ]
MAXIMUM PAY	      [ 	 ] |---------------------------------------
     (VAC. AND S.L.  RATES ARE UNITS EARNED PER PAY INTERVAL)
   VACATION	RATE   [	 ] ACCUMULATED [	 ] USED [	  ]
   SICK LEAVE	RATE   [	 ] ACCUMULATED [	 ] USED [	  ]
   COMPENSATORY TIME		   ACCUMULATED [	 ] USED [	  ]
.sp5
.li
		     Y O U R	C O M P A N Y	 N A M E
			       PAYROLL SYSTEM
  00/00/00		    PAY TRANSACTION ENTRY		  BATCH: nnn





	   RECORD NUMBER [    ]--------------------------------------|
	     |							     |
	     |	EMPLOYEE NO.  [     ]  employee name		     |
	     |	TRANS CODE    [  ]     description		     |
	     |	UNITS	      [ 	   ]			     |
	     |	DATE	      [        ]			     |
	     |							     |
	     |-------------------------------------------------------|

			   DISPLAY EMPLOYEE NAMES [ ]
.bp
.li
		      Y O U R	C O M P A N Y	N A M E
			       PAYROLL SYSTEM
  00/00/00		      PRINT PAY CHECKS


	       TYPE ""Y"" OR PRESS  RETURN  TO PRINT A TEST FORM
	       TYPE ""N"" WHEN THE FORMS ARE PROPERLY ALIGNED [ ]

	       THE LAST CHECK PRINTED WAS CHECK NUMBER	 nnnnn

	       THE STARTING CHECK NUMBER WILL BE	[     ]

	       THE COMPUTER IS PRINTING CHECK NUMBER	 nnnnn
	       THE CHECK IS PAYABLE TO EMPLOYEE NUMBER	 nnnnn

		    PRESS THE ESCAPE KEY TO STOP
		    PRESS THE RETURN KEY TO CONTINUE   [ ]

	       THE CHECK AMOUNT EXCEEDS THE MAXIMUM ALLOWED.
			 THIS CHECK WILL BE VOIDED.

		       ALL CHECKS HAVE BEEN PRINTED
.sp5
.li
		    Y O U R    C O M P A N Y	N A M E
			      PAYROLL SYSTEM
  00/00/00	      EMPLOYEE MASTER LIST PRINT MENU

       1. ALL EMPLOYEES 		   8. ALL HOURLY EMPLOYEES
       2. ALL WEEKLY EMPLOYEES		   9. ALL SALARIED EMPLOYEES
       3. ALL BI-WEEKLY EMPLOYEES	  10. ALL ACTIVE EMPLOYEES
       4. ALL SEMI-MONTHLY EMPLOYEES	  11. ALL ON-LEAVE EMPLOYEES
       5. ALL MONTHLY EMPLOYEES 	  12. ALL TERMINATED EMPLOYEES
       6. ALL WEEK BASED EMPLOYEES	  13. ALL DELETED EMPLOYEES
       7. ALL MONTH BASED EMPLOYEES	  14. A RANGE OF EMPLOYEES

		  ENTER YOUR SELECTION NUMBER: [  ]

		  RANGE FROM [	   ] TO [     ]

		  FULL INFORMATION OR SINGLE LINE (F OR S) [ ]
.bp
.li
	      Y O U R	 C O M P A N Y	  N  A M E
			 PAYROLL SYSTEM
  00/00/00		PRINT 941 REPORT
		company name				date qtr ended
		trade name
		address 				I.D. number
		address
		city		      state  zipcode
		PRINT NAME AND ADDRESS	[ ]
NUMBER OF EMPLOYEES
TOTAL WAGES AND TIPS
TOTAL INCOME TAX WITHHELD
ADJUSTMENT OF WITHHELD INCOME TAX FROM PREVIOUS QUARTERS    [		   ]
ADJUSTED TOTAL OF INCOME TAX WITHHELD
TAXABLE FICA WAGES PAID 		    TIMES 12.26% =
TAXABLE TIPS REPORTED			    TIMES  6.13% =
TOTAL FICA TAXES
ADJUSTMENT OF FICA TAXES				    [		   ]
ADJUSTED TOTAL OF FICA TAXES
TOTAL TAXES
EARNED INCOME CREDIT
NET TAXES
.sp2
.li
DEPOSIT PER END  00/00/00	 LIABILITY	   DATE 	AMOUNT
OVERPAYMENT FROM PREVIOUS QUARTER...........................[		 ]
 FIRST	 | DAYS 1 THROUGH 7	 n,nnn,nnn.nn	[	 ]  [		 ]
 MONTH	 | DAYS 8 THROUGH 15	 n,nnn,nnn.nn	[	 ]  [		 ]
 OF	 | DAYS 16 THROUGH 22	 n,nnn,nnn.nn	[	 ]  [		 ]
 QUARTER | DAYS 23 THROUGH END	 n,nnn,nnn.nn	[	 ]  [		 ]
   FIRST MONTH TOTAL:		 n,nnn,nnn.nn
 SECOND  | DAYS 1 THROUGH 7	 n,nnn,nnn.nn	[	 ]  [		 ]
 MONTH	 | DAYS 8 THROUGH 15	 n,nnn,nnn.nn	[	 ]  [		 ]
 OF	 | DAYS 16 THROUGH 22	 n,nnn,nnn.nn	[	 ]  [		 ]
 QUARTER | DAYS 23 THROUGH END	 n,nnn,nnn.nn	[	 ]  [		 ]
   SECOND MONTH TOTAL:		 n,nnn,nnn.nn
 THIRD	 | DAYS 1 THROUGH 7	 n,nnn,nnn.nn	[	 ]  [		 ]
 MONTH	 | DAYS 8 THROUGH 15	 n,nnn,nnn.nn	[	 ]  [		 ]
 OF	 | DAYS 16 THROUGH 22	 n,nnn,nnn.nn	[	 ]  [		 ]
 QUARTER | DAYS 23 THROUGH END	 n,nnn,nnn.nn	[	 ]  [		 ]
   THIRD MONTH TOTAL:		 n,nnn,nnn.nn
TOTAL FOR QUARTER:		 n,nnn,nnn.nn		    [		 ]
FINAL DEPOSIT FOR QUARTER:......................[	 ]  [		 ]
TOTAL DEPOSITS FOR QUARTER:.................................
OVERPAYMENT:....		  UNDEPOSITED TAXES DUE:....
.bp
.li
		     Y O U R	C O M P A N Y	 N A M E
			       PAYROLL SYSTEM
 00/00/00		      PRINT W-2 REPORT

 ---------------------------------------------------------------------------
				     |
      1  ALL EMPLOYEES		     |	PRINT STATE AND LOCAL AMOUNTS  [ ]
      2  RANGE OF EMPLOYEES	     |
      3  TERMINATED EMPLOYEES	     |	NUMBER OF FORMS PER EMPLOYEE   [  ]
      4  A SINGLE EMPLOYEE	     |
				     |	NUMBER OF LINES BETWEEN FORMS  [  ]
	 ENTER SELECTION [ ]	     |
				     |	STARTING PRINT COLUMN	       [  ]
 ------------------------------------|--------------------------------------
   STARTING EMPLOYEE NUMBER [	  ]  |
       EMPLOYEE NUMBER	 [     ]     |	      PRINT TEST FORM [ ]
   ENDING EMPLOYEE NUMBER   [	  ]  |
 ---------------------------------------------------------------------------

		  NOW PRINTING EMPLOYEE NUMBER nnnnn
.sp6
.li
		      Y O U R	C O M P A N Y	N A M E
			      PAYROLL SYSTEM
  00/00/00		    HISTORY UPDATE MENU



		  1.  UPDATE EMPLOYEE HISTORY TOTALS
		  2.  UPDATE COMPANY LIABILITIES,VACATION ETC.
		  3.  UPDATE PAYMENTS HISTORY



			 ENTER SELECTION NUMBER: [ ]
.bp
.po3
.ll81
.li
		      Y O U R	C O M P A N Y	 N A M E
				PAYROLL SYSTEM
  00/00/00		    UPDATE EMPLOYEE HISTORY
     EMP NO:[	  ]
			 QTD EARNINGS BY TAX CATAGORIES
INCOME:[	  ]  OTHER:[	      ] NO TAXES:[	    ] FICA:[	      ]
		EIC CREDIT:[	      ]     TIPS:[	    ]  NET:[	      ]
QTD TAXES WITHHELD:    FWT:[	      ]      SWT:[	    ]  LWT:[	      ]
		      FICA:[	      ]      SDI:[	    ]
			    QTD EARNINGS AND DEDUCTIONS
SYS: 1:[	  ] 2:[ 	 ] 3:[		] 4:[	       ] 5:[	      ]
EMP: 1:[	  ] 2:[ 	 ] 3:[		] OTHER DEDUCTIONS:[	      ]
UNITS BY PAY TYPE:  1:[ 	 ] 2:[		] 3:[	       ] 4:[	      ]
			 YTD EARNINGS BY TAX CATAGORIES
INCOME:[	  ]  OTHER:[		NO TAXES:[	    ] FICA:[	      ]
		EIC CREDIT:[	      ]     TIPS:[	    ]  NET:[	      ]
YTD TAXES WITHHELD:    FWT:[	      ]      SWT:[	    ]  LWT:[	      ]
		      FICA:[	      ]      SDI:[	    ]
			    YTD EARNINGS AND DEDUCTIONS
SYS: 1:[	  ] 2:[ 	 ] 3:[		] 4:[	       ] 5:[	      ]
EMP: 1:[	  ] 2:[ 	 ] 3:[		] OTHER DEDUCTIONS:[	      ]
UNITS BY PAY TYPE:  1:[ 	 ] 2:[		] 3:[	       ] 4:[	      ]
.sp2
.li
		       Y O U R	 C O M P A N Y	 N A M E
				PAYROLL SYSTEM
  00/00/00		    EMPLOYEE HISTORY TOTALS
			 QTD EARNINGS BY TAX CATAGORIES
INCOME[ 	   ]OTHER[	      ]NO TAXES[	    ]FICA[	      ]
	       EIC CREDIT[	      ]    TIPS[	    ] NET[	      ]
QTD TAXES WITHHELD:   FWT[	      ]     SWT[	    ] LWT[	      ]
		     FICA[	      ]     SDI[	    ]
			    QTD EARNINGS AND DEDUCTIONS
SYS:1[		  ]2[		 ]3[		]4[	       ]5[	      ]
EMP:1[		  ]2[		 ]3[		]OTHER DEDUCTIONS[	      ]
UNITS BY PAY TYPE: 1[		 ]2[		]3[	       ]4[	      ]

			 YTD EARNINGS BY TAX CATAGORIES
INCOME[ 	   ]OTHER[	      ]NO TAXES[	    ]FICA[	      ]
	       EIC CREDIT[	      ]    TIPS[	    ] NET[	      ]
YTD TAXES WITHHELD:   FWT[	      ]     SWT[	    ] LWT[	      ]
		     FICA[	      ]     SDI[	    ]
			    YTD EARNINGS AND DEDUCTIONS
SYS:1[		  ]2[		 ]3[		]4[	       ]5[	      ]
EMP:1[		  ]2[		 ]3[		]OTHER DEDUCTIONS[	      ]
UNITS BY PAY TYPE: 1[		 ]2[		]3[	       ]4[	      ]
.bp
.li
		      Y O U R	C O M P A N Y	N A M E
				PAYROLL SYSTEM
  00/00/00		  UPDATE COMPANY LIABILITIES

		   COMPANY LIABILITIES(NOT TOTALED FROM HISTORY)
				QUARTER TO DATE
	  FICA:[	    ] FUTA:[		] SUI:[ 	   ]
				  YEAR TO DATE
	  FICA:[	    ] FUTA:[		] SUI:[ 	   ]

			 QUARTER TO DATE(NOT TOTALED)
					EARNED		   TAKEN
	       VACATION TIME:	   [		]     [ 	   ]
	       SICK LEAVE:	   [		]     [ 	   ]
	       COMP TIME:	   [		]     [ 	   ]

		    YEAR TO DATE(TOTALED FROM EMPLOYEE RECORDS)
					EARNED		   TAKEN
	       VACATION TIME:	   [		]     [ 	   ]
	       SICK LEAVE:	   [		]     [ 	   ]
	       COMP TIME:	   [		]     [ 	   ]
.sp3
.li
		     Y O U R   C O M P A N Y   N A M E
				PAYROLL SYSTEM
  00/00/00		   UPDATE PAYMENTS HISTORY
			 FWT TAX PAYMENTS FOR QUARTER
		       1. [	       ] ON [	     ]
		       2. [	       ] ON [	     ]
		       3. [	       ] ON [	     ]
		       4. [	       ] ON [	     ]
		       5. [	       ] ON [	     ]
		       6. [	       ] ON [	     ]
		       7. [	       ] ON [	     ]
		       8. [	       ] ON [	     ]
		       9. [	       ] ON [	     ]
		       10.[	       ] ON [	     ]
		       11.[	       ] ON [	     ]
		       12.[	       ] ON [	     ]
			    QUARTERLY TAX PAYMENTS
  FWT:	1.[	       ] 2.[		] 3.[		 ] 4.[		  ]
  FICA: 1.[	       ] 2.[		] 3.[		 ] 4.[		  ]
  FUTA: 1.[	       ] 2.[		] 3.[		 ] 4.[		  ]
.he
.bp
.mt7
.mb9
.po8
.ll65
.ig
CHAPTER 33
D M S	S E T U P    F O R   N O N - S T A N D A R D   C R T ' S
.sp2
.pp5
To create a new CRT.DEF file, run the CRTDEF program, answering all the
requests that pertain to the CRT you are using.  When the program requests
the name of the CRT definition file to create, respond with CRT.  The
program will create a file named CRT.DEF.  This file should always reside
on the
disk which is on the currently logged drive.
All of the other DEF files can be omitted from your work disks or system
master (but not from the original Applewood Computers distribution disk!),
as well as the CRTDEF.INT
program. The CRTDEF program should only have to be executed once to define
the CRT you are using.
.he Appendix: Chapter 33
.pp
The utility program CRTDEF is used to define the physical characteristics of
a CRT display device (e.g., number of rows and columns, screen clear
sequence, cursor addressing sequence, etc.) to the Display Management System.
The program will prompt for
values for all CRT control sequences which are supported by the DMS system.
.pp
The CRTDEF program does not use screen formatting and will work on any serial
CRT.  It will request data for almost every conceivable type of CRT control
procedure.  Not all of these procedures are need by the Payroll
System programs.  The only control sequences that are necessary are
the direct cursor addressing control characters and translation tables,
destructive backspace, and the clear screen character.	Any other requests
can be ignored by pressing RETURN.
Some control items will cause the system to execute faster because it can take
advantage of more sophisticated hardware features.
.pp
If you are modifying an existing CRT definition file, the CRTDEF program
will present a menu so you can choose which section of the program to run by
entering the number of the desired selection.  If you are creating a CRT
defintion file for the first time, the menu does not appear; all of the
requests that are not applicable can be answered with a RETURN.
.pp
When the program issues a request, the current value of the control sequence
is displayed and your are prompted for a new value or acceptance
of the current value.	To accept the current value, just press the RETURN
key.
.pp
When the program starts, it requests the name of the CRT definition
file.  For the Payroll System, always
enter CRT.  This is the name of the definition file that is created by
CRTDEF and the DEF file the Payroll System will look for.
.pp
CRTDEF then requests whether this run will create the initial version
of the file or edit an existing copy.  The first time you run this
program enter a "C" for create.  If you are modifying a DEF file that was
entered incorrectly, enter an "M" for modify.
.cp8
.sp2
THE CONTROL SEQUENCE REQUESTS:
.pp
Most of the CRT control sequences are entered in the form:
.sp
.li
	<number of items>,<item-1>,<item-2>, . . .
.sp
The values of the items are the data characters which will be transmitted
to the CRT device to effect the specified control operation.  These
values are specified as decimal integers between 0-255.  To enter an
ESCAPE, for example, use the decimal equivalent of the ASCII value which
is 27 in decimal.  If a control sequence consisted of an ESCAPE followed by
a left parenthesis, you would enter:
.sp
.li
	2,27,40
.sp
Where the 2 indicates that there are two characters that must be sent to the
CRT for this function, and the 27 and 40 are the decimal representations
of an ESCAPE and a left parenthesis.
.pp
Please note that the CRTDEF program does not check the values for validity.
Valid values range from 0 to 255 but any integer value may be supplied.  The
value modulo 256 will be used.
.pp
The values specified are stored in character strings and are written
to the disk file as such.  CBASIC and CP/M impose restrictions on the
values which may be contained in these strings.  The restrictions are:
.sp
.li
	CP/M:  CTRL-Z (26) is used as an end-of-file marker.
.sp
.li
	CBASIC:  A RETURN (10) or a LINE-FEED (13) character
		 denotes an end of line sequence.
.sp
The way around the restrictions for most terminals which might happen to use
these characters for CRT control is to specify the value as <value>+128.
This sets the ASCII parity bit to ON.  For example, to send a decimal
10, send 138 instead.
.pp
The ASCII terminals will ignore the parity bit because the defined
ASCII set is 7 bits.  (Normally the CBIOS [operating system] will force
the parity bit to a fixed value before transmitting it to the terminal,
so the problem never really arises.)
.sp2
REQUESTS ISSUED BY CRTDEF:
.sp
NUMBER OF ROWS FOR DISPLAY:
.in8
This is the number of rows (lines) of vertical display on your CRT.
.qi
.sp
.li
NUMBER OF COLUMNS FOR CRT DISPLAY:
.in8
This is the number of columns (characters across) displayed on your CRT.
.qi
.sp
.li
SPECIFY CURSOR ROW ADDRESS POSITIONING CONTROL CHARACTERS:
	row 1  current value:  0
		   new value:
.in8
Enter the value required by the CRT to position the cursor at the desired row.
Don't include the position control character sequence, merely the number
that will be translated into a row address.  This request will be issued
once for every possible row.
.qi
.sp
.li
SPECIFY CURSOR COLUMN ADDRESS POSITIONING CONTROL CHARACTERS:
	column 1   current value:  0
		       new value:
.in8
Enter the value required by the CRT to position the cursor at the desired
column.  Don't include the position control character sequence, merely the
number that will be translated into a column address.  This request will
be issued once for every possible column.
.qi
.sp
.li
CRT CONTROL:  SET CURSOR ADDRESS:
.in8
As described above, this request starts with "CRT control:" and thus requires
the number of elements to precede the values of the elements and for them to
be separated by commas.
.sp
Enter the values only if your CRT uses one control character sequence to send
the cursor to both row and column.  Enter the characters required to send the
cursor to a ROW and a COLUMN at the same time.
.qi
.sp
.li
SET ROW ADDRESS PREFIX:
.in8
If the CRT requires different control characters for ROW and COLUMN, enter
the row characters here (the previous request can be left a zero).  If the
CRT address row and column at the same time this can be left zero.
.qi
.sp
.li
SET COLUMN ADDRESS PREFIX:
.in8
If the CRT requires different control characters for ROW and COLUMN, enter
the column characters here.  If the CRT addresses row and column at the same
time this can be left zero.
.qi
.sp
.li
SPECIFY ORDER OF PRESENTATION OF ROW AND COLUMN ADDRESS:
.in8
Enter an "R" or a "C" for the order in which the control characters are to
be sent.
.qi
.sp
.li
SPECIFY FORMAT OF ROW AND COLUMN CURSOR CONTROL CHARACTERS:
.in8
Most CRT's require control information to be sent as hexadecimal characters,
but some require ASCII numbers.  Enter an "H" or a "D" for the correct format.
.qi
.sp
.li
CLEAR SCREEN:
.in8
Enter the characters for clearing the CRT screen.  This is a required entry.
.qi
.sp
.li
CLEAR FOREGROUND DATA FIELDS:
.in8
If your CRT supports different display intensities, enter the
foreground/background
control characters for this and the next 2 requests.
.qi
.sp
.li
ENTER FOREGROUND MODE:

ENTER BACKGROUND MODE:

RETURN CURSOR TO HOME POSITION:

MOVE CURSOR UP ONE ROW:

MOVE CURSOR DOWN ONE ROW:

MOVE CURSOR LEFT ONE COLUMN:

MOVE CURSOR RIGHT ONE COLUMN:

.cp7
.li
DESTRUCTIVE BACKSPACE:
.in8
This sequence must be entered.	This is the sequence of characters
that will erase the previously typed in character.  This will
almost always be a backspace followed by a blank followed by a backspace.
In decimal, this translates to an 8, a 32, another 8.  Key in:
3,8,32,8.
.qi
.sp
.li
SOUND CONSOLE ALARM
.in8
This is normally a 1,7.
.qi
.sp
.li
ERASE TO END OF CURRENT LINE:

ERASE TO END OF SCREEN:

INSERT A NEW LINE INTO THE DISPLAY:

DELETE CURRENT LINE FROM DISPLAY:

NUMBER OF STORAGE LOCATIONS TO BE SET AT CRT INITIALIZATION:
.in8
Specify number of storage locations to be set or press RETURN.	If your
CRT requires memory to be initialized before operating, enter the number
of bytes that must be set at this request.
.qi
.sp
.li
CURRENT STORAGE ENTRY AND VALUES:   0	 0
SET STORAGE ENTRY 1:
.in8
Specify memory location and value.  This request will be
issued for each location
specified in the previous request.  Enter the decimal value for each byte.
.qi
.sp
.li
INITIALIZATION SUBROUTINE ADDRESS:
.in8
If a machine code routine resides in memory that must be branched to on
startup, enter its address in decimal.
.qi
.sp
.li
INITIALIZATION I/O PORT:
.in8
If data needs to be transmitted to an I/O port, enter its number.
.qi
.sp
.li
DATA TO BE TRANSMITTED TO I/O PORT:
.in8
Enter the number of values and the decimal values of the characters that
must be sent to the I/O port.
.qi
.sp
.li
DATA TO BE TRANSMITTED TO CONSOLE:
.in8
If data must be sent to the console device itself at startup time, enter
the number of bytes and their decimal values here.
.qi
.sp
.pp
The requests which follow, or
options 8 and 10 on the CRTDEF menu, are provided to allow you to change
the normal definition of the special function control keys in the system.
If, for example, your operating system or computer does not allow free use
of a CONTROL-A character (enter adding mode), the CRTDEF program would
allow you to define a CONTROL-Y entered at the keyboard to be interpreted
by the program as though you had entered a CONTROL-A.  Because the
facility is completely generalized, it should be possible to compensate
for the special requirements of all of the CP/M derivative operating
systems.
.pp
In addition, the use of control keys can be completely bypassed if
necessary.  A prefixing character can be defined that substitutes for
a control key.	In this manner, for example, an at-sign (@) or an
up arrow (^) can be designated as the prefix character.  Then an
@A sequence  would be interpreted the same as a CONTROL-A.
.pp
The menu option #8 requests
allow a control character translation table to be entered.  It
issues up to 31 requests for each control character possible.  If, for
example, your system has forced you to not use a CONTROL-A (ASCII value
of 1) and you will substitute a CONTROL-Y (ASCII value of 25), when the
request:
.sp
.li
	CHARACTER  25	    CURRENT TRANSLATION   25
	ENTER NEW VALUE (OR STOP) >
.sp
is issued, respond with the value 1.  Don't forget that 8 is the backspace
and 13 is the carriage return.
.pp
The menu option #10 requests are
used only if you cannot use any control characters at all.
They will define instead a prefix character and the individual values
represented by the prefix character followed by one of the 127 possible
ASCII values.
.pp
The first request issued is for the value of the prefix character itself.
Enter its decimal ASCII value.	An at-sign, for example, would be 64.
Remember that once a character is defined as the prefixing character it
cannot be used as a valid data character anywhere else.
This is very important, so choose carefully.
.pp
The next request is used merely to quickly fill all 127 elements of the
table with a default.	The default should be some invalid character
combination such as a zero.
.pp
The program will then issue up to 127 requests for translation of individual
characters.  If you wish for the character sequence "prefix-A" (65), don't
forget to enter the same interpreted value at the lower-case value of "A" (97).
.pp
If you use the prefix character facility of option 10, you should use the
control character translation facility of option 8 to "mask" out all special
function control characters.  This involves setting them to some harmless
but invalid value, such as a zero or 10 (LINE FEED).  All programs will
then reject this character with an explicit error message.
.pp
After all requests have been answered (when the menu is not used) or after
each menu funtion (when the menu is used), the program will display:
.sp
.li
	Specify S(ave), E(dit), or Q(uit)
.sp
If all the data was entered correctly, enter an S.  If the data needs
to be edited, enter an E and the menu will appear, allowing the
data just entered to be modified.  If you are unsure about the entries
appropriate for your device, see your dealer.
.he
.bp
.ig
CHAPTER 34
E R R O R    M E S S A G E S
.sp2
.ce
CBASIC ERROR MESSAGES
.pp5
Applewood Computers software is designed to trap all errors that can arise under
normal circumstances, but due to language limitations and/or hardware
problems, some errors may occur which the programs themselves cannot
handle.  These errors are CBASIC errors.  They appear in the following
form:
.he Appendix: Chapter 34
.sp
.li
		ERROR aa
.sp
where "aa" are two alphabetic characters issued by the CBASIC runtime
package CRUN2.	The occurance of a CBASIC error message usually
indicates a serious problem in your INT or data files.
.sp
.in10
.ti-10
ERROR RE: This is caused by a record in a file which has too few
or too many file delimiter characters (quotes or commas), or
whose length is incorrect.  When a program attempts to
read the record, what is on the disk does not match the record description,
and CRUN2 issues an ERROR RE message.
.sp
This problem may result from hardware
problems like a bad write, a voltage surge, etc., but by far
the most common cause is a bad RAM location, which either changes
some character into a quote or a comma, or changes one of these
file delimters into some other character.  This error is virtually never
the result of a program bug.  Once a bad record is in the file, it must be
repaired or removed with the editor, or the entire file must
be rebuilt.  Bad records can be detected by TYPE'ing the
file and manually searching for the missing or extra
quote or comma.  If you are not familiar with the use of the CP/M editor,
read the CP/M documentation before attempting to change the file.
Also read the File Definition in the Payroll System Specifications
Manual, because some files may require that the header record be
changed if a record is manually deleted.
.sp
.ti-10
ERROR IV: This error occurs if you try to run a program compiled
by CBASIC version 1 with CRUN2, or if the INT file that you are
trying to run is null.	If the latter is the case, recopy the
file from the original disk, using the verify option [vo] of PIP (for
example, PIP A:=B:*.INT[vo]).
.sp
.ti-10
ERROR DZ: This error occurs if you try to divide by zero, or if you try
to run a program compiled under CBASIC2 with CRUN version 1.  Since
none of the Payroll System programs allow division with a
denominator of zero, make sure your runtime package is
indeed CRUN2 (ideally version number 2.04).
.sp
.ti-10
ERROR OM or GARBAGE ON SCREEN DURING AN ENTRY PROGRAM: This means
that the program has run out of memory--should not occur if you have
at least 54K of RAM and a standard-size version of CP/M.
.sp
.ti-10
ERROR SB: Your program files (INT files) may be damaged.  Make new
copies using the CP/M PIP program [vo] option.	For example,
.sp
.li
	  PIP A:=B:*.INT[vo]
.sp
.ti-10
ERROR RG: See ERROR SB.
.qi
.sp2
.ce
PAYROLL SYSTEM ERROR MESSAGES
.sp
.ce
SYSTEMWIDE MESSAGES
.sp
.in6
.ti-6
PY016 CONTROL CHARACTER NOT ACCEPTED
.sp
.ti-6
PYO17 LENGTH LIMIT EXCEEDED
.sp
.ti-6
PY044 INVALID COMMAND SEQUENCE: Contact your distributor.
.sp
.ti-6
PY045 INVALID COMMAND SEQUENCE: Contact your distributor.
.qi
.sp2
.ce
PAY TRANSACTION ENTRY (HRSENT.INT)
.sp
.in6
.ti-6
PY051 EMPLOYEE FILE NOT FOUND ON DRIVE d: Where "d" is a drive name.
Make sure the proper disk has
been inserted.	The employee master file should be named PYEMP.101.
.sp
.ti-6
PY052 EMPLOYEE FILE IS NULL:  Use the CP/M TYPE command to examine the file.
If giving the TYPE command returns the CP/M prompt, erase the null employee
file and bring your backup disks up to date, or re-run the
employee Entry program.
.sp
.ti-6
PY053 Y OR N ONLY
.sp
.ti-6
PY054 INVALID RESPONSE
.sp
.ti-6
PY055 HOURS FILE OF TYPE 000 FOUND:  A previous run the the Pay
Transaction Entry
program ended prematurely.  TYPE the file to determine if it is null.  If
it is, erase it and re-enter the transactions.	If it is not, rename the
file to PYHRS.001, print a proof of the file, and make sure the
transactions are complete and correct.
.sp
.ti-6
PY056 MUST BE NUMERIC INTEGER
.sp
.ti-6
PY057 NOT NUMERIC
.sp
.ti-6
PY058 EXCEEDS NUMBER OF RECORDS ON HOURS FILE
.sp
.ti-6
PY059 INVALID CONTROL CHARACTER: See the documentation.
.sp
.ti-6
PY060 INVALID CHECK DIGIT: The check digit is the rightmost digit of the
five digit employee number.  The check digit is calculated from the other
four digits to insure against miskeys and transposition errors.  A list of
employees and their employee numbers can be printed at any time.
.sp
.ti-6
PY061 INVALID EMPLOYEE NUMBER:	This employee number is not on file.
The short version of the Employee Master List lists the valid
employee numbers.
.sp
.ti-6
PY062 TRANS CODE CANNOT BE LESS THAN ZERO OR GREATER THAN 50:  See Chapter 13
or 26 for a list of valid transaction codes.
.sp
.ti-6
PY063 INVALID DATE: See Chapter 18 for instructions on How to Entry a Date.
.sp
.ti-6
PY064 PR2 FILE NOT FOUND ON SYSTEM DRIVE: Should not occur. Contact your
distributor.
.sp
.ti-6
PY065 INVALID TRANS CODE: See Chapter 13 or 26 for a list of valid transaction
codes.
.sp
.ti-6
PY066 THERE ARE NO HOURS FILE RECORDS:	Make sure the wrong disk was not
inserted, or give the CRTL-A command to begin adding records.
.qi
.sp2
.ce
PAYROLL SYSTEM MENU  (PY.INT)
.sp
.in6
.ti-6
PY101 ONE OF THE TWO PARAMETER ENTRY OPTIONS MUST BE SELECTED:	One of your
parameter files is not present.  At first time startup this is normal.
Make sure the correct disk is inserted.  If the missing file is the PR1
file, it can be recreated by running Parameter Entry.  If it is the
PR2 file, the problem is more serious.	Parameter Entry will create this
file, but the information on it is updated in the course of Payroll System
processing and will most likely be inaccurate.	To bring it up to date may
require the use of an editor.  Check the file specifications to find which
fields should be changed.  This recovery procedure is recommend for those
with substantial data processing knowledge only.
.sp
.ti-6
PY102 DMS ERROR: The CRT.DEF file is invalid for some reason.  Determine
the reason, then PIP a valid CRT.DEF file onto the disk using the [vo]
option of PIP (e.g., PIP A:=B:CRT.DEF[vo]).
.sp
.ti-6
PY103 ERROR DURING CHAIN TO SORT: Check to see if your disk or directory
is full.  If not, contact your distributor.
.sp
.ti-6
PY104 ERROR DURING CHAIN TO SORT: See PY103.
.sp
.ti-6
PY105 THE PROGRAM IS NOT ON THE CURRENTLY LOGGED DISK:	INT files are
expected to be on the currently logged drive.  Make sure the correct
programs disk is inserted.
.sp
.ti-6
PY106 THE SORT PARAMETER FILE IS NOT ON THE SYSTEM DISK:  The Sort Parameter
file PYHOURS.SRT is expected to be on the drive named as the System File
Drive at first time startup.  A new file can be created by running the
Create Sort Parameters program.
.sp
.ti-6
PY107 UNABLE TO OPEN PARAMETER FILE: This error should not occur.  Contact
your distributor.
.sp
.ti-6
PY110 INVALID SECOND PARAMETER FILE:  Invalid information is on the second
parameter file (PYPR2).  A new one can be created by running the Parameter
Entry program, but since this file contains information that is updated
in the course of normal payroll processing, the new file will most likely
be inaccurate.	Check the file specifications and use an editor to change
the file where necessary.  This recovery procedure is recommended only for
those with substantial data processing knowledge.
.sp
.ti-6
PY111 QSORT IS NOT ON THE CURRENTLY LOGGED DISK: Make sure the correct
disk is inserted.
.sp
.ti-6
PY112 THAT'S NOT A VALID RESPONSE: The number of a menu item, ESCAPE, and
CTRL-R are the only valid responses.
.sp
.ti-6
PY114 THE SYSTEM DRIVE HAS NOT BEEN INITIALIZED: This error should not
occur.	Contact your distributor.
.sp
.ti-6
PY115 INVALID DATE: See Chapter 18 for instructions on how to enter a date.
.sp
.ti-6
PY116 INVALID PARAMETER FILE:  The first parameter file (PYPR1) contains
invalid information.  A new file may be created by running Parameter Entry.
.sp
.ti-6
PY117 OPERATOR REQUESTED RETURN TO MENU:  Should not occur.  Contact your
distributor.
.sp
.ti-6
PY118 **ERROR**  THE PROGRAM WAS TERMINATED BEFORE COMPLETION:	Should not
occur.	Contact your distributor.
.sp
.ti-6
PY119 THE PARAMETER ENTRY PROGRAM IS NOT ON THE PROGRAM DISK: Check to make
sure the correct disk has been inserted.
.qi
.sp2
.ce
EMPLOYEE ENTRY (EMPENT.INT)
.sp
.in6
.ti-6
PY131 THIS PROGRAM MODULE MUST BE SELECTED FROM THE PAYROLL MENU:
Payroll System programs must be run from the menu.  Call up the menu
by typing CRUN2 PY.
.sp
.ti-6
PY132 DOUBLE QUOTES (") ARE INVALID CHARACTERS: Use single quotes (') if
needed.
.sp
.ti-6
PY133 ONLY INTEGERS ARE ACCEPTED AT THIS FIELD: Integers are whole numbers
with no decimals (5, but not 5.23).
.sp
.ti-6
PY134 THAT'S NOT AN ACCEPTABLE NUMERIC FORMAT FOR THIS FIELD: See Chapter
18 for instructions on How to Enter a Number.
.sp
.ti-6
PY135 VALID DATES ARE THE ONLY ACCEPTABLE ENTRIES AT THIS FIELD:  See Chapter
18 on How to Enter a Date.
.sp
.ti-6
PY136 Y  OR  N	ARE THE ONLY ACCEPTABLE ENTRIES AT THIS FIELD
.sp
.ti-6
PY138 TYPE  RUN  TO CONTINUE, PRESS THE ESCAPE KEY TO STOP: RUN creates
the Employee Master File, ESCAPE ends the Employee Entry program.
.sp
.ti-6
PY139 SYSTEM STARTUP DETECTED: Informational only.
.sp
.ti-6
PY140 THERE IS NO EMPLOYEE FILE, BUT A HISTORY FILE EXISTS:  The Employee
Master file (PYEMP.101) and the Employee History File (PYHIST.101) are created
together and always have the same number of records.  Make sure your
files are in the drives specified in Parameter Entry.  Bring your backup
files up to date if necessary.
.sp
.ti-6
PY141 THE EMPLOYEE FILE IS INVALID OR IS A NULL FILE: TYPE the file
to see if it is null (if it is, the CP/M prompt simply returns).  Bring
your backup files up to date.
.sp
.ti-6
PY142 UNEXPECTED END OF FILE ON EMPLOYEE FILE: The record count on the
header record does not match the actual number of records.  This error
should not occur.  Bring your backup files up to date to recover.
.sp
.ti-6
PY145 UNABLE TO CREATE THE EMPLOYEE FILE - DISK WRITE-PROTECTED?: Should
not occur.  Contact your distributor.
.sp
.ti-6
PY150 THERE IS NO HISTORY FILE, BUT AN EMPLOYEE FILE EXISTS: Make sure the
correct disks are inserted.  Bring your backups up to date to recover.
.sp
.ti-6
PY151 THE HISTORY FILE IS INVALID OR IS A NULL FILE: TYPE the file to
determine if it is null (the CP/M prompt will return immediately).  Bring your
backups up to date to recover.
.sp
.ti-6
PY152 UNABLE TO ADD THE HISTORY RECORD TO THE HISTORY FILE: See
if your disk is full by the STAT command (for example, A>STAT B:).  If it is
not, contact your distributor.
.sp
.ti-6
PY155 UNABLE TO CREATE THE HISTORY FILE - DISK WRITE-PROTECTED?: Check to
see if the disk is write protected.  If it is not, contact your distributor.
.sp
.ti-6
PY160 UNABLE TO UPDATE THE SECOND PARAMETER FILE: This error should not
occur.	Run the Employee Entry program again.  If this error occurs repeatedly,
contact your distributor.
.sp
.ti-6
PY161 ERROR DURING FILE SAVE: The SAVE did not work properly.  Run Employee
Entry to see if any data has been lost and re-enter it if necessary.  If the
error occurs repeatedly, contact your distributor.
.sp
.ti-6
PY169 THE SECOND PARAMETER FILE DOES NOT EXIST OR IS NULL: This error is
not likely to occur.  Re-run the menu program (CRUN2 PY).  See if
the parameter files (PYPR1.101 and PYPR2.101) are on the correct disk.
.sp
.ti-6
PY171 THAT'S NOT A VALID CONTROL RESPONSE: See the documentation.
.sp
.ti-6
PY172 THERE AREN'T THAT MANY EMPLOYEES ON THE SYSTEM
.sp
.ti-6
PY173 THAT'S NOT A VALID EMPLOYEE NUMBER: Your response contained or
non-numeric character or an invalid check digit.
.sp
.ti-6
PY174 SOCIAL SECURITY NUMBERS MUST BE OF THE FORMAT 000-00-0000: You are
required not type in the dashes (-).
.sp
.ti-6
PY175 THE ONLY THING I KNOW ABOUT SEX IS   M  OR  F
.sp
.ti-6
PY176 EMPLOYEE STATUS CAN BE  A(ACTIVE) L(ON LEAVE) T(TERMINATED): This field
accepts A, L, or T only.
.sp
.ti-6
PY180 THAT'S AN UNUSED EMPLOYEE NUMBER: Employee status has been changed
from Terminated to Deleted by the Close for the Year program.  This number
will be reassigned after lower unused employee numbers are assigned.
.qi
.sp2
.ce
EMPLOYEE RATE ENTRY (RATENT.INT)
.sp
.in6
.ti-6
PY201 THIS PROGRAM MODULE MUST BE SELECTED FROM THE PAYROLL MENU: See PY131.
.sp
.ti-6
PY202 THAT'S NOT A VALID CONTROL RESPONSE: See the documentation.
.sp
.ti-6
PY203 THAT'S NOT A VALID CONTROL RESPONSE: See the documentation.
.sp
.ti-6
PY205 ONLY INTEGERS ARE ACCEPTED AT THIS FIELD: See PY133.
.sp
.ti-6
PY206 THAT'S NOT AN ACCEPTABLE NUMERIC FORMAT FOR THIS FIELD: See PY134.
.sp
.ti-6
PY207 DOUBLE QUOTES (") ARE INVALID CHARACTERS: See PY132.
.sp
.ti-6
PY208 Y  AND  N   ARE THE ONLY ACCEPTABLE ENTRIES AT THIS FIELD
.sp
.ti-6
PY209 THAT'S NOT A VALID DATE: See Chapter 18 on How to Enter a Date.
.sp
.ti-6
PY210 THE EMPLOYEE FILE WAS NOT FOUND: Make sure the correct disks are
inserted and the files on the drives set for them during Parameter Entry.
Bring your backup files up to date to recover.
.sp
.ti-6
PY211 THE EMPLOYEE FILE WAS NULL OR INVALID: See PY141.
.sp
.ti-6
PY212 UNABLE TO READ THE EMPLOYEE RECORD -- END OF FILE FOUND: The counter
in the header record disagrees with the actual number of records on the file.
Bring your backup files up to date.
.sp
.ti-6
PY213 UNABLE TO WRITE THE EMPLOYEE RECORD: Check to see if disk is write
protected.  If it is not, contact your distributor.
.sp
.ti-6
PY216 THAT'S NOT A VALID EMPLOYEE NUMBER:  See PY173.
.sp
.ti-6
PY217 THERE AREN'T THAT MANY EMPLOYEES ON THE SYSTEM
.sp
.ti-6
PY218 THAT'S AN UNUSED EMPLOYEE NUMBER: See PY180.
.sp
.ti-6
PY219 THERE ARE NO EMPLOYEES ON THE SYSTEM: No employees were added
during Employee Entry.	See Chapter 23 for instructions on adding employees.
.sp
.ti-6
PY221 H AND S  ARE THE ONLY VALID ENTRIES AT THIS FIELD: H=Hourly, S=Salaried.
.sp
.ti-6
PY222 THAT'S NOT ONE OF THE PAY INTERVALS CURRENTLY USED:  The pay intervals
to be used are set through the Parameter Entry program.
.sp
.ti-6
PY223 M  AND  S  ARE THE ONLY VALID ENTRIES AT THIS FIELD: M=Married, S=Single.
.sp
.ti-6
PY226 THERE ARE FOUR TAX EXCLUSION TYPES NUMBERED ONE THRU FOUR: See Chapter
10 or Chapter 22 for the four tax exclusion codes.
.sp
.ti-6
PY229 THIS EMPLOYEE HAS BEEN EXEMPTED FROM FEDERAL TAX WITHHOLDING: An
employee's exemption status can be changed through the Employee Ded/Earn
and Cost program.
.sp
.ti-6
PY230 THIS EMPLOYEE HAS BEEN EXEMPTED FROM STATE TAX WITHHOLDING: See PY229.
.sp
.ti-6
PY231 THIS EMPLOYEE HAS BEEN EXEMPTED FROM LOCAL TAX WITHHOLDING: See PY229.
.qi
.sp2
.ce
EMPLOYEE EARN/DED AND COST (EMPDENT.INT)
.sp
.in6
.ti-6
PY241 THIS PROGRAM MODULE MUST BE SELECTED FROM THE PAYROLL MENU: See PY131.
.sp
.ti-6
PY242 THAT'S NOT A VALID CONTROL RESPONSE: See the documentation.
.sp
.ti-6
PY243 THAT'S NOT A VALID CONTROL RESPONSE: See the documentation.
.sp
.ti-6
PY245 ONLY INTEGERS ARE ACCEPTED AT THIS FIELD: See PY133.
.sp
.ti-6
PY246 THAT'S NOT AN ACCEPTABLE NUMERIC FORMAT FOR THIS FIELD: See PY134.
.sp
.ti-6
PY247 DOUBLE QUOTES (") ARE INVALID CHARACTERS: See PY132.
.sp
.ti-6
PY248 Y  AND  N   ARE THE ONLY ACCEPTABLE ENTRIES AT THIS FIELD
.sp
.ti-6
PY249 THAT'S NOT A VALID DATE: See Chapter 18 on How to Enter a Date.
.sp
.ti-6
PY250 THE EMPLOYEE FILE WAS NOT FOUND: See PY210.
.sp
.ti-6
PY251 THE EMPLOYEE FILE WAS NULL OR INVALID: See PY141.
.sp
.ti-6
PY252 UNABLE TO READ THE EMPLOYEE RECORD -- END OF FILE FOUND: See PY212.
.sp
.ti-6
PY253 UNABLE TO WRITE THE EMPLOYEE RECORD: See PY213.
.sp
.ti-6
PY256 THAT'S NOT A VALID EMPLOYEE NUMBER: See PY173.
.sp
.ti-6
PY257 THERE AREN'T THAT MANY EMPLOYEES ON THE SYSTEM
.sp
.ti-6
PY258 THAT'S AN UNUSED EMPLOYEE NUMBER: See PY180.
.sp
.ti-6
PY259 THERE ARE NO EMPLOYEES ON THE SYSTEM: See PY219.
.sp
.ti-6
PY264 AN INVALID CHECK CATEGORY HAS BEEN ENTERED: See Chapter 15 and
Chapter 28.
.sp
.ti-6
PY265 E AND D  ARE THE ONLY VALID ENTRIES AT THIS FIELD: E=Earning
D=Deduction.
.sp
.ti-6
PY266 THERE ARE FOUR TAX EXCLUSION TYPES NUMBERED ONE THRU FOUR: See PY226.
.sp
.ti-6
PY267 A AND P  ARE THE ONLY VALID ENTRIES AT THIS FIELD: A=flat Amount,
P=Percentage of gross taxable earnings.
.sp
.ti-6
PY268 CHECK CATEGORIES CAN BE  2-7 (EARNINGS) OR 9-16 (DEDUCTIONS):
See Chapter 15 and Chapter 28 for a treatment of check categories.
.sp
.ti-6
PY269 THAT'S NOT A VALID PAYROLL SYSTEM ACCOUNT:  Accounts must always
be referred to be their Payroll System Account Reference Number, which
is the relative position of the account on the account file.  An account
must have been entered onto the Account file through the Account Entry
program before it is considered a valid account.
.sp
.ti-6
PY270 THE ACCOUNT DISTRIBUTION PERCENTAGES MUST TOTAL 100.0%: The program
will not let you leave until the percentages total 100.0.
.qi
.sp2
.ce
PARAMETER ENTRY (PYPARENT.INT)
.sp
.in6
.ti-6
PY281 PROGRAM MUST BE SELECTED FROM THE PAYROLL MENU: See PY131.
.sp
.ti-6
PY281 ONLY INTEGERS ARE ACCEPTED AT THIS FIELD: See PY133.
.sp
.ti-6
PY283 THAT'S NOT AN ACCEPTABLE NUMERIC FORMAT FOR THIS FIELD: See PY134.
.sp
.ti-6
PY284 DOUBLE QUOTES ARE INVALID CHARACTERS: See PY132.
.sp
.ti-6
PY285 Y  AND  N  ARE THE ONLY ACCEPTABLE ENTRIES AT THIS FIELD
.sp
.ti-6
PY286 THAT'S NOT A VALID DATE: See Chapter 18 for instructions on How
to Enter a Date.
.sp
.ti-6
PY287 THAT'S NOT AN ACCEPTABLE ENTRY FOR THIS FIELD
.sp
.ti-6
PY288 THAT'S NOT IN THE ACCEPTABLE RANGE FOR THIS FIELD
.sp
.ti-6
PY289 THAT'S NOT A VALID PAYROLL SYSTEM ACCOUNT: See PY269.
.sp
.ti-6
PY290  THE ACCOUNT ENTRY MODULE IS NOT ON THE PROGRAM DISK:  Your system
may be set up incorrectly. See Chapter 19 for detailed instructions.
Run Account Entry to create the default account file.
.sp
.ti-6
PY291 THAT PARAMETER ENTRY MODULE IS NOT ON THE PROGRAM DISK:  Your system
may be set up incorrectly.  See Chapter 19 for detailed instructions.
.sp
.ti-6
PY292 UNABLE TO UPDATE PARAMETER FILE AT END OF PROGRAM: Make sure the
parameter files are on the correct disk.  Rerun Parameter Entry to create
a new PYPR1 file. This error should not occur.
.sp
.ti-6
PY293 UNABLE TO UPDATE SECOND PARAMETER FILE AT END OF PROGRAM:  Make sure
the PYPR2 file is on the correct disk.	This error should not occur.
.sp
.ti-6
PY294 THAT'S NOT A VALID PAY INTERVAL: Valid pay intervals are W(eekly),
B(i-weekly), S(emi-monthly), and M(onthly).
.sp
.ti-6
PY295 COULDN'T CREATE THE PARAMETER FILE:  Check to see if your disk
if write-protected or the directory full.
.sp
.ti-6
PY296 COULDN'T WRITE TO THE PARAMETER FILE: Check to see if your disk is
full.
.sp
.ti-6
PY297 COULDN'T CREATE THE SECOND PARAMETER FILE: See PY295.
.sp
.ti-6
PY298 COULDN'T WRITE TO THE SECOND PARAMETER FILE: See PY296.
.sp
.ti-6
PY300 THE SYSTEM DRIVE MUST BE	A  OR  B: Other drives are not allowed.
.qi
.sp2
.ce
PARAMETER ENTRY (PYPAR2.INT)
.sp
.in6
.ti-6
PY311 PROGRAM MUST BE SELECTED FROM THE PARAMETER ENTRY MENU: To call
up the Parameter Entry menu, call up the Payroll System by CRUN2 PY,
and then select the Parameter Entry program from the system menu.
.sp
.ti-6
PY311 ONLY INTEGERS ARE ACCEPTED AT THIS FIELD: See PY133.
.sp
.ti-6
PY313 THAT'S NOT AN ACCEPTABLE NUMERIC FORMAT FOR THIS FIELD: See PY134.
.sp
.ti-6
PY314 DOUBLE QUOTES ARE INVALID CHARACTERS: See PY132.
.sp
.ti-6
PY315 Y  AND  N  ARE THE ONLY ACCEPTABLE ENTRIES AT THIS FIELD
.sp
.ti-6
PY316 THAT'S NOT A VALID DATE: See Chapter 18 for instructions on How
to Enter a Date.
.sp
.ti-6
PY317 THAT'S NOT AN ACCEPTABLE ENTRY FOR THIS FIELD
.sp
.ti-6
PY318 THAT'S NOT IN THE ACCEPTABLE RANGE FOR THIS FIELD
.sp
.ti-6
PY319 THAT'S NOT A VALID PAYROLL SYSTEM ACCOUNT: See PY269.
.sp
.ti-6
PY322 UNABLE TO UPDATE PARAMETER FILE AT END OF PROGRAM: See PY292.
.sp
.ti-6
PY323 UNABLE TO UPDATE SECOND PARAMETER FILE AT END OF PROGRAM: See PY293.
.sp
.ti-6
PY324 THAT'S NOT A VALID PAY INTERVAL: See PY294.
.qi
.sp2
.ce
PARAMETER ENTRY (PYPAR3.INT)
.sp
.in6
.ti-6
PY331 PROGRAM MUST BE SELECTED FROM THE PARAMETER ENTRY MENU: See PY311.
.sp
.ti-6
PY332 ONLY INTEGERS ARE ACCEPTED AT THIS FIELD: See PY131.
.sp
.ti-6
PY333 THAT'S NOT AN ACCEPTABLE NUMERIC FORMAT FOR THIS FIELD: See PY134.
.sp
.ti-6
PY334 DOUBLE QUOTES ARE INVALID CHARACTERS: See PY132.
.sp
.ti-6
PY335 Y  AND  N  ARE THE ONLY ACCEPTABLE ENTRIES AT THIS FIELD
.sp
.ti-6
PY336 THAT'S NOT A VALID DATE: See Chapter 18 for instructions on How
to Enter a Date.
.sp
.ti-6
PY337 THAT'S NOT AN ACCEPTABLE ENTRY FOR THIS FIELD
.sp
.ti-6
PY338 THAT'S NOT IN THE ACCEPTABLE RANGE FOR THIS FIELD
.sp
.ti-6
PY339 THAT'S NOT A VALID PAYROLL SYSTEM ACCOUNT: See PY269.
.sp
.ti-6
PY342 UNABLE TO UPDATE PARAMETER FILE AT END OF PROGRAM: See PY292.
.sp
.ti-6
PY344 THAT'S NOT A VALID DRIVE: Valid drive names under CP/M are A, B,
C, and D.
.qi
.sp2
.ce
SYSTEM DEDUCTION ENTRY MENU (DEDENT.INT)
.sp
.in6
.ti-6
PY351 DEDENT MUST BE RUN FROM PY MENU: To call up the Payroll System menu,
type CRUN2 PY.
.sp
.ti-6
PY353 PYDED - DEDUCTION FILE NOT PRESENT: At first time startup
the program should ask if you want
to create it.
.sp
.ti-6
PY354 IMPROPERLY FORMED NUMBER OR OUT OF RANGE: The numbers 1 through
3 are the only valid responses.
.qi
.sp2
.ce
STANDARD DEDUCTION RATES (DEDENT1.INT)
.sp
.in6
.ti-6
PY351 DEDENT MUST BE RUN FROM PY MENU: To call up the Payroll System menu,
type CRUN2 PY.
.sp
.ti-6
PY352 INVALID RESPONSE
.sp
.ti-6
PY353 PYDED - DEDUCTION FILE NOT PRESENT: Have you switched disks?  This
error should not occur under normal circumstances.  Recreate the file through
the System Deduction Entry program.
.sp
.ti-6
PY354 IMPROPERLY FORMED NUMBER OR OUT OF RANGE: See Chapter 18 for How to
Enter a Number.  The maximum rate must be less than 100.0%.
.sp
.ti-6
PY355 ANSWER MUST BE 'Y' OR 'N'
.qi
.sp2
.ce
FEDERAL WITHOLDING TABLES (DEDENT2.INT)
.sp
.in6
.ti-6
PY351 DEDENT MUST BE RUN FROM PY MENU: To call up the Payroll System menu,
type CRUN2 PY.
.sp
.ti-6
PY352 INVALID RESPONSE
.sp
.ti-6
PY353 PYDED - DEDUCTION FILE NOT PRESENT: See PY353 under Standard Deduction
Rates above.
.sp
.ti-6
PY354 IMPROPERLY FORMED NUMBER OR OUT OF RANGE: See PY354 under Standard
Deduction Rates above.
.sp
.ti-6
PY356 FWT CUTOFFS AND PERCENTS MUST BE IN ASCENDING ORDER:  Change
the incorrect data. See Chapter 10 for a discussion of how to enter
Federal Tax Withholding information.
.sp
.ti-6
PY357 INVALID USE OF CONTROL CHARACTER: See the documentation.
.qi
.sp2
.ce
SYSTEM WIDE EARNINGS AND DEDUCTIONS (DEDENT3.INT)
.sp
.in6
.ti-6
PY351 DEDENT MUST BE RUN FROM PY MENU: To call up the Payroll System
menu, type CRUN2 PY.
.sp
.ti-6
PY352 INVALID RESPONSE
.sp
.ti-6
PY353 PYDED - DEDUCTION FILE NOT PRESENT: See PY353 under Standard Deduction
Rates above.
.sp
.ti-6
PY354 IMPROPERLY FORMED NUMBER OR OUT OF RANGE
.sp
.ti-6
PY355 ANSWER MUST BE 'Y' OR 'N'
.sp
.ti-6
PY358 ANSWER MUST BE 'E' OR 'D': E=Earning D=Deduction.
.sp
.ti-6
PY359 ANSWER MUST BE 'A' OR 'P': A=flat Amount P=Percentage of gross
taxable pay.
.sp
.ti-6
PY360 CHK CAT MUST BE 2-7 FOR EARNINGS, 9-16 FOR DEDUCTIONS: See Chapter 15
and Chapter 28.
.qi
.sp2
.ce
ACCOUNT ENTRY (PYACTENT.INT)
.sp
.in6
.ti-6
PY451 INVALID RESPONSE
.sp
.ti-6
PY452 PROGRAM NOT RUN FROM MENU:  This program must be selected for the
system menu.  To call up the Payroll System menu, type CRUN2 PY.
.sp
.ti-6
PY454 UNEXPECTED END OF ACCOUNT FILE: The counter on the header record
disagrees with the actual number of records on the file.  Bring your
backup file up to date, or recreate the account file.
.sp
.ti-6
PY455 RECORD INFORMATION IS NOT VALID: Enter a valid account number
into the record or give the CTRL-X command to cancel the addition or
change.
.sp
.ti-6
PY456 INVALID RECORD NUMBER: Your entry was out of the range of account
records, or not an integer.
.sp
.ti-6
PY459 CHART OF ACCOUNTS FILE NOT FOUND,INSERT CORRECT DISK: If you have
chosen the General Ledger interface option, this message prompts you to
insert the GL data disk with the Chart of Accounts file into the drive
you set as the GL Date Drive during parameter entry.  Be sure the
two character file suffix that identifies the Chart of Accounts file as
that of your company has been correctly entered in Parameter Entry.
.sp
.ti-6
PY460 UNEXPECTED END OF CHART OF ACCOUNTS FILE: This error should not
occur.	Contact your distributor.
.sp
.ti-6
PY461 INVALID ACCOUNT NUMBER: If the GL interface option has been chosen,
the account number you enter must be a number already on the Chart of
Accounts.  A four digit Chart of Accounts number is valid.  However,
you may not enter an account number with a zero in the fourth position
(1120 or 112034) since this signifies a title account.	If GL interface
is not selected, any 1 to 6 digit number is acceptable.
.sp
.ti-6
PY462 MISSING PR2 FILE, NOT UPDATED: Did you reinsert the disk with the
second parameter file?	May occur if you try to give the CTRL-S command
to save newly entered or changed information while the disk with the
parameter file is switched out to make room for the GL data disk.  If
so, end the program, reinsert the disk, and then Save the file.
.sp
.ti-6
PY463 ACCOUNT FILE NOT PRESENT,INSERT CORRECT DISK: Is the Account file
on the disk you removed to make room for the GL data disk?  Find or
recreate the file.
.qi
.sp2
.ce
CREATE SORT PARAMETER FILE (PYSRTPRM.INT)
.sp
.in6
.ti-6
PY482 UNABLE TO CREATE PYHOURS.SRT:  Check to see if there is space on
your disk or if the disk is write protected.  If there is space and
the disk is write-enabled, contact your distributor since this error
should not occur.
.qi
.sp2
.ce
APPLY PAYROLL (APPLY1.INT)
.sp
.in6
.ti-6
PY501 INVALID DATE: See Chapter 18 for instructions on How to Enter
a Date.
.sp
.ti-6
PY502 CHECK DATE IS NOT IN CURRENT QUARTER: If you want the checks to
carry that pay date, you'll have to run the end of period processing
steps first.  See Chapters 16, 29, and 30.
.sp
.ti-6
PY503 CHECK DATE IS NOT IN CURRENT YEAR: If you want the checks to carry
that pay date, you'll have to close for the fourth quarter and close
for the year first.  See Chapters 16, 29, and 30.
.sp
.ti-6
PY504 MUST BE M OR W: M=Month-based employees, W=Week based employees.
.sp
.ti-6
PY505 CURRENT CHECKS NOT PRINTED: The Print Pay Checks and Print Pay
Check Register programs must be run before another payroll application
may take place.  Pay Check printing is not required if you did not order
automatic check writing during parameter entry.
.sp
.ti-6
PY506 CHECK DATE EARLIER THAN PREVIOUS CHECK DATE:  The check date must
always be past the last check date.
.sp
.ti-6
PY507 PYPAY.100 FOUND: The previous payroll application ended prematurely.
Determine the reason, then rerun the apply using updated backup files.
.sp
.ti-6
PY508 CURRENT PAYROLL JOURNAL NOT PRINTED: You must print at least one
copy of the Payroll Journal report before you can run the Apply Payroll
program again.
.sp
.ti-6
PY509 EMP FILE NOT FOUND: Make sure the correct disk has been inserted,
and the file is on the drive specified for it in Parameter Entry.
If necessary, bring your backup files up to date and use them.
.sp
.ti-6
PY510 DED FILE NOT FOUND: See PY509.
.sp
.ti-6
PY511 SWT OR CALx FILE NOT FOUND: See PY509.
.sp
.ti-6
PY512 LWT FILE NOT FOUND: See PY509.
.sp
.ti-6
PY513 HIS FILE NOT FOUND: See PY509.
.sp
.ti-6
PY514 YES, ESCAPE OR BACK ONLY: The only valid responses to this request
are YES to continue processing, ESCAPE to end the program, and CTRL-B
to go back to the previous request to change your response.
.sp
.ti-6
PY515 ACT FILE NOT FOUND: See PY509.
.sp
.ti-6
PY516 PR2 FILE NOT FOUND: See PY509.
.sp
.ti-6
PY517 YES OR ESCAPE ONLY: Type YES to continue, and press ESCAPE to end the
program.
.sp
.ti-6
PY518 YES OR ESCAPE ONLY: The Company History file does not exist until
after the first payroll application.  The first time you run the Apply
Payroll program, type YES to continue.	If the file has been created by
a previous application, end the current application with ESCAPE, and
check to see if the correct disk has been inserted or if the Company
History file is on the drive you set for it during Parameter Entry.
.sp
.ti-6
PY519 CAN'T APPLY.  QUARTER NEEDS ENDING: You're trying to apply too
many times in one quarter.
.sp
.ti-6
PY520 CAN'T APPLY.  QUARTER NEEDS ENDING: See PY519.
.sp
.ti-6
PY521 CAN'T APPLY.  YEAR NEEDS ENDING: You've closed for the fourth
quarter, but not for the year.
.sp
.ti-6
PY522 CURRENT REGISTER NOT PRINTED: You must print the Check Register
before the program allows another payroll application.
.sp
.ti-6
PY523 CHK.100 FILE FOUND: The previous payroll application ended prematurely.
Determine the reason, then bring your backup files up to date and
rerun the application.
.qi
.sp2
.ce
APPLY PAYROLL (APPLY2.INT)
.sp
.in6
.ti-6
PY552 NOT PROPERLY CHAINED: This error should not occur. Contact your
distributor.
.sp
.ti-6
PY553 HOURS TRANS REJECTED FOR TERMINATED OR ON-LEAVE EMP:  The Apply
Payroll program found pay transactions for terminated or on leave employees.
If you want the transaction to be put through, change the employee's status
through the Employee Entry program.
.sp2
.ce
APPLY PAYROLL  (APPLY3.INT)
.sp
.in6
.ti-6
PY562 NOT PROPERLY CHAINED: This error should not occur.  Contact your
distributor.
.sp
.ti-6
PY563 NET PAY EXCEEDS EMPLOYEE MAXIMUM: Informational message. You can
change the MAX PAY amount for any employee through the Employee Rate
Entry program.
.sp
.ti-6
PY564 NET PAY EXCEEDS SYSTEM LIMIT - CHECK WILL BE VOIDED: Informational
message. You can suppress voiding or change maximum amount through the
Parameter Entry program.
.sp
.ti-6
PY565 NET PAY CALCULATED FOR LESS THAN ZERO - NO CHECK MADE: Make sure
the employee's earning, deduction, and pay transaction information is
correct.
.sp2
.ce
PRINT 940 REPORT (PRINT940.INT)
.sp
.in6
.ti-6
PY581 PR2 FILE NOT FOUND: This error should not occur.	Contact your
distributor.
.sp
.ti-6
PY582 COH FILE NOT FOUND: Check to make sure the correct disks were inserted
and that the file is on the drive specified during Parameter Entry.
.sp
.ti-6
PY583 DED FILE NOT FOUND: See PY582.
.sp
.ti-6
PY584 SYSTEM HAS NOT ACHIEVED YEAR-END STATUS: The program allows you to
continue, but the report will be inaccurate.
.sp
.ti-6
PY585 UNABLE TO WRITE TO PR2: This error should not occur.  Contact your
distributor.
.sp
.ti-6
PY586 HIS FILE NOT FOUND: See PY582.
.sp
.ti-6
PY587 UNEXPECTED EOF ON HIS FILE: The counter in the header record disagrees
with the actual number of records on the Employee History file.  Try to
determine the reason before you continue.  Bring your backup files up to
recover.
.qi
.sp2
.ce
PRINT 941 REPORT (PRINT941.INT)
.sp
.in6
.ti-6
PY601 SYSTEM HAS NOT ACHIEVED QUARTER-END STATUS: The program allows you to
continue, but the report generated will be inaccurate.
.sp
.ti-6
PY602 PR2 FILE NOT FOUND: This error should not occur.	Contact your
distributor.
.sp
.ti-6
PY603 COH FILE NOT FOUND: See PY582.
.sp
.ti-6
PY604 UNABLE TO WRITE PR2 FILE: This error should not occur.  Contact your
distributor.
.sp
.ti-6
PY605 'Y' OR 'N' ONLY
.sp
.ti-6
PY606 NOT NUMERIC
.sp
.ti-6
PY607 INVALID DATE: See Chapter 18 for How to Enter a Date.
.sp
.ti-6
PY609 INVALID RESPONSE
.sp
.ti-6
PY610 PR2 FILE NOT FOUND--SECTION TWO: This error should not occur.  Contact
your distributor.
.sp
.ti-6
PY611 UNABLE TO WRITE PR2--SECTION TWO:  See PY610.
.sp
.ti-6
PY612 COH FILE NOT FOUND--SECTION TWO: See PY610.
.qi
.sp2
.ce
PRINT W2 REPORT (PRINTW2.INT)
.sp
.in6
.ti-6
PY621 SYSTEM HAS NOT ACHIEVED YEAR-END STATE:  Informational message.  W-2
forms are normally produced at the end of the year, although they may be
printed at any time.
.sp
.ti-6
PY622 EMPLOYEE FILE NOT FOUND: See PY582.
.sp
.ti-6
PY623 UNEXPECTED EOF ON EMPLOYEE FILE:	The count of employee records held
on the PR2 file disagrees with the actual number.  Determine the reason
for this, then bring your backup files up to date to recover.
.sp
.ti-6
PY624 HISTORY FILE NOT FOUND: See PY582.
.sp
.ti-6
PY625 UNEXPECTED EOF ON HISTORY FILE:  The count of employee records held
on the PR2 file disagrees with the actual number of records on the
Employee History file.	Determine the reason for this, then bring your
backup files up to date to recover.
.sp
.ti-6
PY626 PR2 FILE NOT FOUND: This error should not occur.	Contact your
distributor.
.sp
.ti-6
PY627 UNABLE TO WRITE TO PR2 FILE
.sp
.ti-6
PY628 INVALID RESPONSE
.sp
.ti-6
PY629 Y, YES, N, OR NO ONLY
.sp
.ti-6
PY630 NUMERIC INTEGER ONLY
.sp
.ti-6
PY631 INVALID CHECK DIGIT
.sp
.ti-6
PY632 EMPLOYEE NUMBER TOO HIGH
.sp
.ti-6
PY633 STARTING EMPLOYEE NUMBER CANNOT BE GREATER THAN ENDING NUMBER
.sp
.ti-6
PY634 OUT OF RANGE
.qi
.sp2
.ce
PRINT PAY CHECKS (PYCHECKS)
.sp
.in6
.ti-6
PY641 THIS MODULE MUST BE SELECTED FROM THE PAYROLL SYSTEM MENU: See PY131.
.sp
.ti-6
PY642 THE SECOND PARAMETER FILE IS NOT ON THE SYSTEM DRIVE: Make certain
you have the correct disk inserted.
.sp
.ti-6
PY643 CHECK FILE WAS NOT RENAMED FOLLOWING EARLIER PROCESSING:	A check
file was left with the filetype of 100. Try to
determine why the rename failed.  TYPE the file to see if it is null
(the CP/M prompt will simply return).  If it is not null, rename it
to 101 and continue.
.sp
.ti-6
PY644 NO CHECK FILE FOUND: Check to make sure the correct disks are inserted.
The file should have been created by the last run of the Apply Payroll program.
.sp
.ti-6
PY645 CHECK FILE RENAME WAS UNSUCESSFUL:  A rename failed.  Your disk
may be write-protected; check for this, then try again.  If it fails
repeatedly, contact your distributor.
.sp
.ti-6
PY646 THE CHECK FILE IS A NULL FILE: Some problem must have occurred during
earlier processing.  Use your backup files to rerun the last payroll
application.
.sp
.ti-6
PY647 UNABLE TO READ CHECK RECORD:  The counter in the header record disagrees
with the actual number or records on the file.	Use your backup files to
rerun the last payroll application.
.sp
.ti-6
PY648 ERROR WHILE WRITING CHECK RECORD: Your disk may be write-protected.
If the Check file has the wrong filetype of 100, rename to type 101.
.sp
.ti-6
PY649 NO EMPLOYEE MASTER FILE FOUND: Check to make sure the correct disk
is inserted.  Bring your backup files up to date and use them.
.sp
.ti-6
PY650 THE EMPLOYEE FILE IS A NULL FILE: Bring your backup files up to date
and use them.
.sp
.ti-6
PY651 UNABLE TO READ EMPLOYEE MASTER RECORD:  The count of records that is
kept on the header record is wrong.  Check to see if data
is missing from the employee file.  Bring your backup files up to date
and use them.
.sp
.ti-6
PY652 UNABLE TO UPDATE THE SECOND PARAMETER FILE:  This error should not
occur.	Contact your distributor.
.sp
.ti-6
PY660 THAT'S NOT A VALID RESPONSE: Valid responses are Y or RETURN to
print another test form, N to begin printing checks, and ESCAPE to end
the program.
.sp
.ti-6
PY661 THAT'S NOT A VALID CHECK NUMBER: Enter a valid 5 digit check number.
The number you enter at this request will be the number of the first
check printed by the computer.
.sp
.ti-6
PY662 ESC AND RETURN ARE THE ONLY RESPONSES ACCEPTED AT THIS FIELD
.qi
.sp2
.ce
PRINT EMPLOYEE MASTER LIST (EMPPRINT.INT)
.sp
.in6
.ti-6
PY671 INVALID RESPONSE
.sp
.ti-6
PY672 PROGRAM NOT RUN FROM MENU: See PY131.
.sp
.ti-6
PY674 UNEXPECTED END OF EMPLOYEE FILE: If you have not yet created the
Employee Master file through the Employee Entry program, do so now.
This error should not occur if a valid Employee Master file exists.
.sp
.ti-6
PY675 SELECTION MUST BE 1 TO 14
.sp
.ti-6
PY676 INVALID OR NON-EXISTANT EMPLOYEE NUMBER: The Employee Master List
shows all the valid employee numbers.
.sp
.ti-6
PY677 FROM NUMBER MUST NOT BE GREATER THAN TO NUMBER
.sp
.ti-6
PY678 TO NUMBER MUST NOT BE LESS THAN FROM NUMBER
.sp
.ti-6
PY679 ANSWER MUST BE 'F' OR 'S': F=Full detail, S=Short version.
.sp
.ti-6
PY680 UNEXPECTED END OF ACCOUNT FILE: The Account file is not present, or
the counter in the header record disagrees with the actual number of
records on the file.  Check to see if the correct disk is inserted,
build a new account file, or use your backup file.
.sp
.ti-6
PY681 INVALID ACCOUNT NUMBER: An account number set for an employee's
GL Cost Distribution or an Employee Specific deduction is invalid.
A three digit Payroll System Account Reference numbers is the only valid
response to an account number request.
.qi
.sp2
.ce
PRINT PAY CHECK REGISTER (PYRGSTR.INT)
.sp
.in6
.ti-6
PY701 INVALID RESPONSE
.sp
.ti-6
PY702 PROGRAM NOT RUN FROM MENU: See PY131.
.sp
.ti-6
PY704 UNEXPECTED END OF EMPLOYEE FILE: The Employee Master file may not be
present.  Make sure you inserted the correct disk.  If the Employee Master
file is present, contact your distributor.
.sp
.ti-6
PY705 UNEXPECTED END OF CHECK FILE: The check file may not be present.
Determine whether the last payroll application was run to successful
completion.  If this is not the problem, contact your distributor.
.qi
.sp2
.ce
CLOSE FOR QUARTER, CLOSE FOR YEAR (PYPEREND.INT)
.sp
.in6
.ti-6
PY721 EMP FILE NOT FOUND: Check to make sure the correct disks have
been inserted, and that the file is on the drive you specified for it
during Parameter Entry.
.sp
.ti-6
PY722 HIS FILE NOT FOUND: See PY721.
.sp
.ti-6
PY723 COH FILE NOT FOUND: See PY721.
.sp
.ti-6
PY724 UNEXPECTED EOF ON COH FILE: This error should not occur. Contact
your distributor.
.sp
.ti-6
PY725 PR2 FILE NOT FOUND: See PY721.
.sp
.ti-6
PY726 CAN'T END YEAR.  FOURTH QUARTER NOT CLOSED: Close for the fourth
quarter before closing for the year.  See Chapters 16, 29, and 30.
.sp
.ti-6
PY727 CAN'T END YEAR.  FOURTH QUARTER NOT CLOSED: See PY726.
.sp
.ti-6
PY728 CAN'T END YEAR.  940 NOT PRINTED: The 940 report must be printed
before you can run the Close for the Year program.  See Chapter 30.
.sp
.ti-6
PY729 CAN'T END YEAR.  W2 NOT PRINTED:  W-2 forms must be printed
for all employees (Active, On Leave, and Terminated) before you can
run the Close for the Year program.  See Chapter 30.
.sp
.ti-6
PY730 CAN'T END QUARTER.  NOT ENOUGH WEEK BASED APPLIES:  There must
be at least 13 weekly applications and at least 7 bi-weekly application before
you can close for the quarter.
.sp
.ti-6
PY731 CAN'T END QUARTER.  NOT ENOUGH MONTH BASED APPLIES: There must
be 3 monthly applications, and 6 semi-monthly applications before
you can close for the quarter.
.sp
.ti-6
PY732 CAN'T END QUARTER.  941 NOT PRINTED: Must print the 941 Report
before you close for the quarter.  See Chapter 16 and Chapter 29.
.sp
.ti-6
PY733 CAN'T END YEAR.  JUST ENDED YEAR:  You've already run the Close
for the Year program.
.qi
.sp2
.ce
PRINT PAYROLL JOURNAL (PYJOURN.INT)
.sp
.in6
.ti-6
PY761 PAY FILE NOT FOUND: See PY721.
.sp
.ti-6
PY762 EMP FILE NOT FOUND: See PY721.
.sp
.ti-6
PY763 ACT FILE NOT FOUND: See PY721.
.sp
.ti-6
PY764 DED FILE NOT FOUND: See PY721.
.sp
.ti-6
PY765 NET PAY EXCEEDS SYSTEM MAX. CHECK WILL BE VOIDED: See PY564.
.sp
.ti-6
PY767 NET PAY LESS THAN ZERO. CHECK WILL BE VOIDED: See PY564.
.qi
.sp2
.ce
PRINT PAYROLL JOURNAL (PYJOURN1.INT)
.sp
.in6
.ti-6
PY771 ACT FILE NOT FOUND: This error should not occur.	Contact your
distributor.
.sp
.ti-6
PY772 UNEXPECTED EOF ON PAY FILE: The file exists but has no header record.
This error should not occur. Contact your distributor.
.sp
.ti-6
PY773 BATCH TOTAL NOT ZERO: If you are not using the GL interface option
this error is not serious.  If you are using GL interface, the batch of
GL transactions
created by the last payroll application is incorrect.  This error means
the accounts have not been set up correctly. For example, you may
have given the number of a cash account for accrued liabilities.  In general,
the Ledger Account Summary printed at the end of the Payroll Journal will
be off by twice the offending amount.  Change the incorrect account information
and rerun the application with updated
backup files, or change the batch through the General Ledger facilities and
correct the accounts before the next application.  See Chapter 9 and Chapter
21.
.qi
.sp2
.ce
TRANSACTION PROOF (HRSPROOF.INT)
.sp
.in6
.ti-6
PY781 INVALID RESPONSE
.sp
.ti-6
PY782 PROGRAM NOT RUN FROM MENU
.sp
.ti-6
PY784 UNEXPECTED END OF TRANSACTION FILE: The Pay Transaction (Hours) file
may not be present.  Make sure it has been created.  If it is present,
contact your distributor.
.sp
.ti-6
PY785 UNEXPECTED END OF EMPLOYEE FILE:	The Employee Master file may not
be present.  Check to make sure the correct disk has been inserted.  If
the file is present, contact your distributor.
.qi
.sp2
.ce
PRINT HISTORY REPORT (HISPRINT.INT)
.sp
.in6
.ti-6
PY801 INVALID RESPONSE
.sp
.ti-6
PY802 PROGRAM NOT RUN FROM MENU: See PY131.
.sp
.ti-6
PY803 INVALID USE OF CONTROL CHARACTER
.sp
.ti-6
PY804 UNEXPECTED END OF HISTORY FILE: The count of employee history records
held on the header record disagrees with the actual number.
Determine the reason for this, then bring your backup files up to date to
recover.
.sp
.ti-6
PY805 HISTORY FILE NOT FOUND: See PY582.
.sp
.ti-6
PY806 EMPLOYEE FILE NOT FOUND: See PY582.
.sp
.ti-6
PY807 UNEXPECTED END OF EMPLOYEE FILE: There are fewer records on the
Employee Master file than on the Employee History file.  Determine the
reason, then bring your backup files up to date to recover.
.sp
.ti-6
PY808 DEDUCTION FILE NOT FOUND: See PY582.
.qi
.sp2
.ce
PRINT VACATION REPORT  (VACPRINT.INT)
.sp
.in6
.ti-6
PY821 INVALID RESPONSE
.sp
.ti-6
PY822 PROGRAM NOT RUN FROM MENU: See PY131.
.sp
.ti-6
PY823 INVALID USE OF CONTROL CHARACTER
.sp
.ti-6
PY824 UNEXPECTED END OF EMPLOYEE FILE: See PY623.
.sp
.ti-6
PY825 EMPLOYEE FILE NOT FOUND: See PY582.
.qi
.sp2
.ce
HISTORY UPDATE (PYUPDHIS.INT)
.sp
.in6
.ti-6
PY921 INVALID RESPONSE
.sp
.ti-6
PY922 PROGRAM NOT RUN FROM MENU: See PY131.
.sp
.ti-6
PY923 INVALID USE OF CONTROL CHARACTER: The CTRL-A command is not allowed.
An employee must first have been entered through the Employee Entry program.
.sp
.ti-6
PY924 UNEXPECTED END OF FILE: Make sure the Employee Master file, the
Employee History file, and the Company History files are on the correct
disk.  If this is not the problem, contact your distributor.
.sp
.ti-6
PY925 NON-NUMERIC INPUT OR EXCEEDS 999,999.99
.sp
.ti-6
PY929 HISTORY UPDATE MAY NOT RUN AFTER APPLY HAS RUN:  The History Update
program not be run after the first run of the Apply Payroll program in order to
protect the auditability of the Payroll System.
.sp
.ti-6
PY930 EMPLOYEE FILE AND HISTORY FILE MUST BE PRESENT: Check to see if the
correct disks have been inserted or that these files have been created
through Employee Entry.
.sp
.ti-6
PY931 INVALID OR NON-EXISTANT EMPLOYEE NUMBER
.sp
.ti-6
PY933 NON-NUMERIC INPUT OR EXCEEDS 9,999,999.99
.sp
.ti-6
PY934 IMPROPERLY FORMED DATE: See Chapter 18 for instructions on How
to Enter a Date.
.qi
.sp2
.ce
PRINT ACCOUNTS	(PYACTPRT.INT)
.sp
.in6
.ti-6
PY961 INVALID RESPONSE
.sp
.ti-6
PY962 PROGRAM NOT RUN FROM MENU: See PY131.
.sp
.ti-6
PY963 INVALID USE OF CONTROL CHARACTER
.sp
.ti-6
PY964 UNEXPECTED END OF ACCOUNT FILE: The counter on the header record
disagrees with the number of records actually on the file.  Determine the
reason, then bring your backup account file up to date or create a new one.
.sp
.ti-6
PY965 ACCOUNT FILE NOT FOUND: See PY582.
.qi
.sp3
.cp12
.ce
QSORT ERROR MESSAGES
.pp
If the QSORT program encounters a problem during processing it will print
a numbered error message to the system console and stop running.  The error
message will contain a descriptive phrase that should usually be enough to
indicate to the user what the problem is.  If more information is desired
the message can be found in the following section by using the error number
that preceeds each error message in the format:
.sp
.li
	QSnn   msg
.sp
where nn is the error number and msg is the descriptive phrase.
.sp
.in5
.ti-5
QS01 MISSING OR INVALID SORT PARAMETERS:
This message indicates one of the following:
.in8
a) The parameter file name was omitted from the command-line when
the QSORT program was initiated.
.br
b) The parameter file name on the command-line was incorrectly
typed, was too long, had embedded blanks, had the wrong drive
specification or the file is not cataloged with the file type
of 'SRT'.
.br
c) The key specification parameters are incorrectly formed.
.sp
.in5
.ti-5
QS02 OPEN ERROR:
This message is issued if the QSORT attempts to open a file that does not
exist on the specified drive.  Check the directory for the name of the
input file.
.sp
.ti-5
QS03 CLOSE ERROR:
This message is issued if the QSORT program attempts to close a file that
is not present.  If this error occurs it may be due to a system error or
to a diskette being changed during the sort process.
.sp
.ti-5
QS04 NO DIR SPACE:
This message indicates that the program is attempting to create a new file
and there is not enough room in the diskette directory.  The file being
created may be the output file or one of the sort work files.  Remove any
unnecessary files from the diskette.
.sp
.ti-5
QS05 INPUT TOO LARGE:
This message will be issued if the QSORT program cannot fit the input
data into the maximum allowable number of work files or if the available
memory is too small for the merging routine.
.sp
.ti-5
QS06 not used
.sp
.ti-5
QS07 OUT OF DISK SPACE:
This message is issued if the sort program runs out of room on a disk write.
This can occur on the output file or one of the sort work files.  Remove any
unnecessary files from the diskette.
.sp
.ti-5
QS08 FILE ERROR:
This message indicates that CP/M cannot locate a directory entry for an output
or work file.  This error only occurs on a write and is an indication of a
possible serious system error.
.sp
.ti-5
QS09 NO DIR SPACE:
This message is similar to QS04 except that it will only occur when a file,
either output or sort work, is being extended by the system.  A file must be
extended every 16K bytes.
.sp
.ti-5
QS10 not used:
.sp
.ti-5
QS11 READ ERROR:
This message indicates a serious program error and should never occur.	It
means that the sort program is attempting to read past the end of file on
an input or work file.
.qi
.he
.bp
.ce
CHAPTER 35
.sp
.ce
E R R O R    R E C O V E R Y
.pp5
When certain programs end abnormally they may leave a file with an
inproper filetype.  Under normal conditions the filetypes of most
files in the system should end with a 1 (one) to indicate the
file is a valid current file.  Since the 1 is changed to a 0 (zero)
when the file is being processed, and renamed back to 1 when
processing is complete, you will never see a zero as the third digit
unless the program ended before the rename took place, or the rename
failed for some other reason.
.he Appendix: Chapter 35
.pp
Before renaming the file with the CP/M REN command, you should try to
learn why the program ended prematurely or the rename failed.  It is
quite possible that the problem is hardware related, and it may even be
due to a problem with the operating system.
.pp
To rename a file, give the REN command in the following form:
.sp
.li
	    REN D:nnnnnnnn.ttt=D:nnnnnnnn.ttt
.sp
where D is a drive reference, nnnnnnnn an 8 character filename, and ttt
the 3 character filetype.  The filename and type to the left of the
equal sign is that to which the file on the right will be renamed.  For
example, to rename PYPAY.100 to PYPAY.101 give the command,
.sp
.li
	    REN B:PYPAY.101=B:PYPAY.100
.sp
if the PYPAY.100 file is on B and you want the PYPAP.101 file on B also.
The rename cannot take place if a file with the same name as that
to which the other file is to be renamed already exists on the
designated drive.
.pp
Certain error conditions may arise because a disk or its directory is full,
or the disk write-protected.  If the disk is write-protected, remove
the write-protect tab or apply a write-enable tab (whichever is
appropriate for your system).  Unless your disk contains many other files
besides those used and created by the Payroll System, the disk directory
should not be full.  However, since the number of files a directory can
hold does vary, you should find out the capacity of the directory
you are using.
.pp
To check to see if a disk is full, use the CP/M STAT command.  Although
this command should be on the currently logged drive, it can give you
the status of files on other drives.  For example, to find the number
of bytes remaining on the B drive with STAT on the A drive, type
the command,
.sp
.li
	    STAT B:
.sp
See the CP/M Features and Facilities manual for complete instructions
on how to use the STAT command.
.pp
If a header record is incorrect, you may be able to correct it using a file
editor.  Check the file specifications for the necessary
information.  Be certain to get the record length correct.  This
recovery procedure is suggested only for those with substantial
operating experience or data processing knowledge.  In any case,
back up your files before attempting to modify them with an editor.
For others, the best solution is to bring your backup files up to date.
.sp2
.ce
USING SDCOPY
.pp
A utility program is distributed with every Payroll System that can be
used to back up single density disks in a minimum amount of time.
The program, called SDCOPY.COM, also verified the copied data to ensure
the copy is identical to the original.
.pp
To execute the program, type SDCOPY at the CP/M prompt and follow with
RETURN.  The program responds with the log-on message and "mount" message
shown in figure 35.1.  The program is now stopped and either or both
disks may be removed from their drives.  The disk to be copied (the "source"
disk) must be
place in drive A and the backup disk in drive B.  The "destination" disk
in drive B may be a freshly formatted, blank disk, or any other disk.
Any data or programs on the destination disk are irretrievably destroyed
and replaced with the contents of the source disk.
.sp
.cp14
.li
    A>sdcopy rs1
    SDCOPY VER:2.0
    PUT SOURCE ON A, DESTINATION ON B, THEN TYPE RETURN  [return]

    SUCCESSFUL COPY
    PUT SOURCE ON A, DESTINATION ON B, THEN TYPE RETURN  [return]

    SUCCESSFULL COPY
    PUT SOURCE ON A, DESTINATION ON B, THEN TYPE RETURN  [ctrl-c]

    INSERT SYSTEM DISKETTE, THEN TYPE RETURN TO REBOOT	 [return]
    A>

	 (Figure 35.1: The SDCOPY mount message)
.pp
When both the source and destination disks have been inserted, hit RETURN
and the SDCOPY program will begin to copy.  When the copy is complete,
you can repeat (if the repeat option has been specified, as above), or
enter CTRL-C to end the program.  After entering CTRL-C, put
the program disk in drive A, and hit RETURN. The system will then be in the
normal operating mode awaiting further commands.  The new backup
copy should be placed in a secure location.
.pp
The SDCOPY program offers several options.  Each option may be specified
by typing a single charcter following the program name and separated by
at least one space.  If an "R" is typed, the copy program repeats.  The
program repeatedly displays the mount message (allowing both disks
to be changed) until a CTRL-C is entered instead of the RETURN key.  This
allows several disks to be conveniently copied at one time.
.pp
If an "S" is typed, the copy program stops copying at the first
empty data "track".  For partially full disks this option can
speed up program execution.  This option may only be used to copy source
disks that have been "SYSGEN'ed".
.pp
Because the access times of various manufacturer's disk controllers vary,
data is spread around the revolving disks at varying distances.  Consecutive
items of data in a file may not necessarily occupy physically consecutive
positions on the disk.	This is termed "skewing".  The SDCOPY program
will run up to ten times faster if the skewing is properly set.  Nine
differenct skews are possible and they are specified by typing a number
from one to nine after the name SDCOPY.  Trial and error is the best
way to find the optimum skewing for each manufacturer's drives.  A
one is the best for Digital System and a four is best for Micromation.
If no skewing number is used, the SDCOPY program defaults to four.
.pp
Any combination of the options may be specified for any run of
SDCOPY.  The characters may be in any order and need not separated by
spaces.
.pp
The SDCOPY program will only run on standard CP/M systems that are based
on 8 inch, 3740 compatible floppy disks.  The program will not work
on double density systems.  The disk drive manufacturer should have
provided a full disk copy program with the hardware.
.bp
.mt7
.mb7
.po8
.ll65
.ig
CHAPTER 36
C O M M A N D	 S U M M A R Y
.sp2
.in8
.ti-8
ESCAPE: Ends the current function or program.
.sp
.ti-8
RETURN: Enters data and advances to next request; accepts default value;
moves cursor forward through DMS screen.
.sp
.ti-8
CTRL-H: Backs cursor up one position, erasing character in that position.
.sp
.ti-8
CTRL-R: Refreshes garbled DMS screen.
.sp
.ti-8
CTRL-X: Cancels current addition or change, works only before ESCAPE has
been pressed.
.sp
.ti-8
CTRL-A: Add records to the file.
.sp
.ti-8
CTRL-D: Delete record currently accessed.
.sp
.ti-8
CTRL-S: Save records added or changed so far.
.sp
.ti-8
CTRL-N: Get the next record on the file.
.sp
.ti-8
CTRL-B: Get the previous record on the file; go back one request.
.qi
.sp3
.po5
.li
	  TRANSACTION CODES			 TAX EXCLUSION CODES

 0  Rate0 Pay	       12  Sick Pay	      1  Subject to all taxes
 1  Rate1 Pay	       13  Vacation Pay 	 and percent deductions
 2  Rate2 Pay	       14  Other Excl. Pay
 3  Rate3 Pay	       15  Exp. Reimburse.    2  Subject to all taxes
 4  Rate4 Pay	       17  Other Pay		 and percent deductions
 5  Vacation Taken     18  Tips Reported	 except FICA
 6  Sick Leave Taken   27  Advance Repay.
 7  Comp Time Taken    28  FWT Add-On	      3  Subject to all taxes
 8  Comp Time Earned   29  SWT Add-On		 and percent deductions
 9  Bonus	       30  LWT Add-On		 except FWT, SWT, & LWT
10  Tips Collected     31  FICA Add-On
11  Advance	       50  Other Deds.	      4  Not subject to taxes or
						 percent deductions.

		TO CALL UP THE PAYROLL SYSTEM:

   CRUN2 PY MM/DD/YY  [date optional, not accepted at first startup]
.br
.po8

